%----------------
% !TEX program = xelatex
% Requires LaTeXbangla 2.x (TeX Live 2024+ / current MiKTeX)

\documentclass[letterpaper,oneside,openany,12pt]{book}
\usepackage[DIV=14,BCOR=2mm,headinclude=true,footinclude=false]{typearea}

% (geometry removed: it conflicted with typearea, which overrode it anyway)

%--------------------------------------------------------------------
% Bangla: LaTeXbangla loads fontspec + polyglossia, makes Bangla the
% default language, switches Bangla/English fonts automatically, and
% uses Bangla digits for counters (chapters, pages, equations, refs).
%--------------------------------------------------------------------
\usepackage{latexbangla}
\usepackage{fancyhdr}
\pagestyle{fancy}
\fancyhf{}  % clear default header and footer

% Chapter name only, without "অধ্যায় ১"
\renewcommand{\chaptermark}[1]{\markboth{#1}{#1}}

% Even pages: book name, odd pages: chapter name
\fancyhead[L]{\ifodd\value{page}\else\thetitle\fi}
  \fancyhead[R]{\ifodd\value{page}\leftmark\fi}
\fancyfoot[C]{\thepage}

\renewcommand{\headrulewidth}{0.4pt} 
% Kalpurush files bundled with the project (kalpurush.ttf, Kalpurush-Italic.ttf)
\setbanglamainfont{kalpurush}[
  Path=./,
  Extension=.ttf,
  UprightFont=kalpurush,
  ItalicFont=Kalpurush-Italic
]
% If Kalpurush is installed system-wide, this one line is enough instead:
% \setbanglamainfont{Kalpurush}

% Optional English font (otherwise Latin Modern):
% \setmainfont{...}

% Old helpers kept as aliases so existing chapters still compile.
% They are no longer needed in new text: just type Bangla directly.
\newcommand{\bn}[1]{#1}
\newcommand{\bnit}[1]{\textit{#1}}

%--------------------------------------------------------------------
% Packages
%--------------------------------------------------------------------
\usepackage{xcolor}
\usepackage{graphicx}
\usepackage{lastpage}
\usepackage{epigraph}
\usepackage{makeidx}
\usepackage{titling}   % control \maketitle formatting
\usepackage{multicol}
\usepackage{subcaption}
\usepackage{xspace}
\usepackage[skins]{tcolorbox}
\usepackage{rotating}
\usepackage{amsmath,amstext,amssymb}
\usepackage{enumitem}
\usepackage[mathscr]{eucal}
\usepackage{aas_macros}
\usepackage{bm}
\usepackage{changepage}
\usepackage{listings}
\usepackage{siunitx}
\usepackage{lipsum}
\usepackage[square, numbers, comma, sort&compress, sectionbib]{natbib}
\usepackage{chapterbib}
\usepackage{quotchap}
\usepackage[backref=section, breaklinks, plainpages=false, colorlinks=true,
  anchorcolor=cyan, linkcolor=cyan, citecolor=cyan, filecolor=magenta,
  urlcolor=magenta, linktoc=all]{hyperref}

%--------------------------------------------------------------------
% Dialogue environments for the characters
%--------------------------------------------------------------------

%--------------------------------------------------------------------
% Units used in the Introduction
%--------------------------------------------------------------------
\DeclareSIUnit{\Msun}{\ensuremath{M_\odot}}
\DeclareSIUnit{\gauss}{G}
\DeclareSIUnit{\parsec}{pc}
\DeclareSIUnit{\yr}{yr}
\DeclareSIUnit{\years}{yr}

%--------------------------------------------------------------------
% Chapter numbers (quotchap): large grey Bangla digits
%--------------------------------------------------------------------
\renewcommand*{\chapnumfont}{\fontsize{100}{130}\selectfont\color{chaptergrey}}

% Grey epigraph
\newcommand{\grayepigraph}[2]{%
  {\color{gray}\epigraph{#1}{#2}}%
}

% vertical spacing between paragraphs
\setlength{\parskip}{0.5em}

\frenchspacing
\tolerance=5000

%--------------------------------------------------------------------
% Title
%--------------------------------------------------------------------
\title{ইতিহাসে‌‌ জ্যোতির্বিজ্ঞান}

\author{  \textbf{লেখক}\\ এন্টন পেনাকোপ\\[5em]
  \textbf{অনুবাদক} \\[1em]আতিকুর রহমান\\মোঃ আব্দুস সামী আকন্দ ‌‌}
\date{}

\pretitle{\centering\fontsize{70}{50}\selectfont}
\posttitle{\par\vskip 1em\normalfont}
\preauthor{\centering\Large}
\postauthor{\par\vskip 1em\normalfont}
\predate{\centering\large}
\postdate{\par\vskip 2em\normalfont}

\begin{document}
\maketitle

\thispagestyle{plain}
\begin{center}
% \Large{Abstract}
\end{center}
\vspace{1.5cm}
% \lipsum[1]

\newpage
% \epigraph{``\textit{বল মহাবিশ্বের মহাকাশ ফাড়ি’\\
% চন্দ্র সূর্য গ্রহ তারা ছাড়ি’\\
% ভূলোক দ্যুলোক গোলক ভেদিয়া\\
% খোদার আসন ‘আরশ’ ছেদিয়া,\\
% উঠিয়াছি চির-বিস্ময় আমি বিশ্ববিধাতৃর}''}{\textit{নজরুল}}
\tableofcontents
\part{প্রাচীনকালের জ্যোতির্বিজ্ঞান}

% Chapter 1
\section{নক্ষত্র ও জীবন }

জ্যোতির্বিজ্ঞানের শুরুটা হয় প্রাগৈতিহাসিক সময়ের আগে থেকে, যে সময়ের কোন রেকর্ড খুজে পাওয়া যায় না।

% ch2.tex

\chapter{চাষাবাদ আর বর্ষপঞ্জি}\label{ch:2}

\epigraph{``\textit{
আমায় প্রশ্ন করে নীল ধ্রুব তারা,\\
আর কত কাল আমি রব দিশেহারা
}''}{\textit{সলিল চৌধুরী}}

% From the multitude of tribes living in a condition of barbarism, large states arose between 4000 and 1000 BC in the fertile plains of China, India, Mesopotamia and Egypt.
খ্রিস্টপূর্ব ৪০০০ থেকে ১০০০ অব্দের মধ্যে চীন, ভারত, মেসোপটেমিয়া ও মিশরের উর্বর সমভূমিতে বসবাসকারী বিভিন্ন জনগোষ্ঠী ধীরে ধীরে গড়ে তোলে বড় বড় রাষ্ট্র ও সভ্যতা।
% They represented a higher stage of culture, of which written documents have been transmitted to posterity.
তারা সংস্কৃতি ও সভ্যতার দিক থেকে আরও উন্নত পর্যায়ে পৌঁছেছিল। তাদের সেই অগ্রগতির প্রমাণ আজও পাওয়া যায় পরবর্তী প্রজন্মের জন্য রেখে যাওয়া বিভিন্ন লিখিত দলিল ও নথিপত্রে।
% The formerly independent peasant communities and town-ships, with their own chieftains and urban kings, their local gods and forms of worship, merged into larger political units.
আগে যেসব কৃষকসমাজ ও নগরকেন্দ্রিক জনপদ নিজেদের মতো স্বাধীনভাবে পরিচালিত হতো, যাদের ছিল নিজস্ব সর্দার বা নগরশাসক, স্থানীয় দেবদেবী ও উপাসনার রীতি, সেগুলো ধীরে ধীরে বৃহত্তর রাজনৈতিক ব্যবস্থার সঙ্গে একীভূত হয়ে যায়।
% The extraordinary fertility of the alluvial silt,  which yielded a surplus product, provided a living for a separate class of rulers and officiate.
উর্বর পলিমাটিতে প্রচুর ফসল ফলত। প্রয়োজন মেটানোর পর যে বাড়তি ফসল থেকে যেত, তা শাসক ও সরকারি কর্মকর্তাদের মতো পৃথক একটি শ্রেণির ভরণপোষণের ব্যবস্থা করত।
% This organization first arose from the need for the centralized regulation of the water.
পানির প্রবাহ ও বণ্টন কেন্দ্রীয়ভাবে নিয়ন্ত্রণের প্রয়োজন থেকেই প্রথম এই প্রশাসনিক ব্যবস্থার উদ্ভব ঘটে।
% The large rivers irrigating these plains the Nile, the Euphrates, the Hwang Ho filled their beds with silt, overflowed in certain months and flooded the fields, devastating or fertilizing them, or at times excavating a new bed.
এই সমভূমিগুলোর ওপর দিয়ে বয়ে যাওয়া নীল নদ, ইউফ্রেতিস ও হোয়াং হোর মতো বড় বড় নদী পলিমাটি জমিয়ে নিজেদের তলদেশ ভরাট করে ফেলত। বছরের নির্দিষ্ট সময়ে নদীগুলো দু'কূল ছাপিয়ে প্লাবিত করত আশপাশের কৃষিজমি। কখনো সেই বন্যায় ফসল ও জমির ক্ষতি হতো, আবার কখনো পলিমাটি ছড়িয়ে জমি হয়ে উঠত আরও উর্বর। কোনো কোনো সময় নদী তার পুরোনো গতিপথ ছেড়ে নতুন পথও তৈরি করত।
% The water had to be directed continuously and checked by dykes, by a deepening of the rivers, or by the digging of canals.
পানির প্রবাহ ঠিক রাখতে এবং তা নিয়ন্ত্রণে রাখতে নিয়মিত বাঁধ নির্মাণ, নদীর তলদেশ গভীর করা কিংবা নতুন খাল খননের প্রয়োজন হতো।
% Such control could not be left to the individual districts with their often conflicting interests. Centralized regulation was necessary, and only a strong central authority could guarantee that local interests would not prevail over general interests. Only then were fertility and prosperity assured.
বিভিন্ন অঞ্চলের স্বার্থ প্রায়ই পরস্পরের সঙ্গে সাংঘর্ষিক হতো। তাই পানির এই ব্যবস্থাপনার দায়িত্ব আলাদা আলাদা অঞ্চলের হাতে ছেড়ে দেওয়া সম্ভব ছিল না। প্রয়োজন ছিল কেন্দ্রীয় নিয়ন্ত্রণের, আর একমাত্র শক্তিশালী কেন্দ্রীয় শাসনব্যবস্থাই নিশ্চিত করতে পারত যে সামগ্রিক স্বার্থের চেয়ে কোনো অঞ্চলের নিজস্ব স্বার্থ বড় হয়ে উঠবে না। এভাবেই নিশ্চিত করা সম্ভব হতো জমির উর্বরতা ও সমৃদ্ধি।
% But when the country fell asunder into petty warring principalities and the dykes and canals were neglected, the soil dried up or was flooded, and the people starved: the wrath of the gods then lay on the land.
কিন্তু দেশ যখন ভেঙে ছোট ছোট বিবদমান রাজ্যে বিভক্ত হয়ে পড়ত এবং সেগুলোর মধ্যে যুদ্ধ লেগেই থাকত, তখন বাঁধ ও খালগুলোর যথাযথ রক্ষণাবেক্ষণ হতো না। ফলে কখনো জমি শুকিয়ে ফসল নষ্ট হতো, আবার কখনো বন্যায় ডুবে যেত। শেষ পর্যন্ত মানুষ দুর্ভিক্ষের মুখে পড়ত। তখন এই বিপর্যয়কে মনে করা হতো দেবতাদের অভিশাপ বা রোষের ফল।



% A strong central power was needed, secondly, for the defence of the fertile plains against the war-like inhabitants of the adjacent mountains or deserts.
শক্তিশালী কেন্দ্রীয় শাসনব্যবস্থার প্রয়োজন ছিল আরও একটি কারণে—পাশের পাহাড়ি ও মরু অঞ্চলে বসবাসকারী যুদ্ধপ্রবণ জনগোষ্ঠীর আক্রমণ থেকে উর্বর সমভূমিগুলোকে রক্ষা করার জন্য।
% These people could find only a meagre living on their own land, so they made it their business to plunder and to force tribute from their prosperous neighbours.
নিজেদের ভূমিতে জীবিকা নির্বাহ করাই ছিল তাদের জন্য কঠিন। তাই তারা আশপাশের সমৃদ্ধ জনপদে লুটতরাজ চালাত এবং সেখানকার মানুষের কাছ থেকে জোর করে কর বা খাজনা আদায় করত।
% A division of tasks was necessary; a caste of warriors developed who, with their chief as king, became a ruling class controlling the surplus product of the farmers.
এভাবে কাজের দায়িত্ব ভাগ করে দেওয়ার প্রয়োজন থেকে একটি পৃথক যোদ্ধাশ্রেণির উদ্ভব ঘটে। তাদের নেতা রাজা হিসেবে প্রতিষ্ঠিত হন এবং এই যোদ্ধাশ্রেণিই ধীরে ধীরে শাসকগোষ্ঠীতে পরিণত হয়ে কৃষকদের উৎপাদিত বাড়তি ফসলের ওপর নিয়ন্ত্রণ প্রতিষ্ঠা করে।
% Or the roving neighbours turned from marauders into conquerors and settled as a ruling aristocracy in the midst of the farmers, protecting them from other aggressors. In either case the result was a strong centralized state power.
কখনো আশপাশের যাযাবর জনগোষ্ঠী লুটপাটের পরিবর্তে ভূখণ্ড দখল করে সেখানে স্থায়ীভাবে বসতি গড়ে তুলত। এভাবে তারা কৃষকদের মাঝে শাসক অভিজাত শ্রেণি হিসেবে প্রতিষ্ঠিত হতো এবং অন্য আক্রমণকারীদের হাত থেকে কৃষকদের রক্ষা করত। উভয় ক্ষেত্রেই শেষ পর্যন্ত গড়ে উঠত একটি শক্তিশালী কেন্দ্রীয় শাসনব্যবস্থা।



% This story, repeated from time to time, represents in brief outline the history of these lands. Over and over again barbarian peoples invaded and subjected them, sometimes remaining as a thin upper layer, like the Manchus in China, the Mongols in India, the Hyksos in Egypt, and sometimes as an entire people mixing with or replacing the former inhabitants, like the Aryans in India and the Semites in Mesopotamia.
এই ভূখণ্ডগুলোর ইতিহাস পর্যালোচনা করলে বারবার একই ঐতিহাসিক ধারার পুনরাবৃত্তি চোখে পড়ে। বিভিন্ন সময়ে নানা যুদ্ধবাজ জাতিগোষ্ঠী এসব ভূখণ্ডে আক্রমণ চালিয়েছে। কখনো তারা বিজিত অঞ্চলের ওপর শাসকগোষ্ঠীর একটি ক্ষুদ্র উপরিস্তর হিসেবে নিজেদের প্রতিষ্ঠিত করেছে, অথচ মূল শাসনক্ষমতা বজায় রেখেছে নিজস্ব মাতৃভূমিতে—যেমন চীনে মাঞ্চুরা, ভারতে মোগলরা এবং মিশরে হিকসসরা। আবার কখনো তারা স্থানীয় জনগোষ্ঠীর সঙ্গে মিশে গিয়ে তাদেরই অংশ হয়ে উঠেছে, যেমন ভারতে আর্যরা এবং মেসোপটেমিয়ায় সেমিটিক জনগোষ্ঠী। কোনো কোনো ক্ষেত্রে স্থানীয় জনগোষ্ঠীকে সম্পূর্ণরূপে বিতাড়িত বা উৎখাতও করেছে।
% Though much culture was destroyed in the conquest, the invaders afterwards adopted and assimilated the existing higher state of civilization, and often imparted fresh vigour to it. After some generations, having lost their savage strength, the conquerors themselves became an object of attack by new aggressors.
আগ্রাসনের ফলে অনেক সময় স্থানীয় সংস্কৃতি ও ঐতিহ্য বিলুপ্ত হলেও, বিজয়ী জনগোষ্ঠী পরবর্তীকালে বিজিত অঞ্চলের অপেক্ষাকৃত উন্নত সভ্যতা ও সংস্কৃতিকে গ্রহণ করে নিজেদের মধ্যে আত্মস্থ করেছে। এর মধ্য দিয়ে তাদের সমাজে নতুন প্রাণশক্তি ও উদ্যমের সঞ্চার হয়েছে। তবে কয়েক প্রজন্মের ব্যবধানে বিজয়ী শাসকগোষ্ঠীও নিজেদের আগের সেই দুর্ধর্ষ শক্তি ও সামরিক দাপট হারিয়ে ফেলত। একসময় তারাই নতুন আক্রমণকারীদের আগ্রাসনের শিকারে পরিণত হতো।




In such empires the prince was the legislator., the chief justice and the head of a government of officials, who as leaders of the civil departments constituted a second ruling class alongside the military. Usually it consisted ofthe priests, who had been the local intellectual leaders, and now were organized into an official hierarchy. The priesthood held the spiritual leadership of state and society in its hands. It held the theoretical and general knowledge needed in the process of production; this was the source of its prestige and its social power. Where agriculture was the main occupation, knowledge of the calendar and of the seasons was their domain. Religion too in those times one with state and society was centralized; the local deities of the chief townships were conjoined into a pantheon under an upper god, whereas the local gods of smaller or conquered sites were degraded to a lower order of spirits or were incorporated into other gods. Thus in ancient Babylonia the goddesses Innina, Nisaba and Nana were united with the later Ishtar.

The fact that Borsippa became a suburb of the rapidly growing city of Babylon found its theological expression in its city god Nabo becoming the son ofMarduk. Whereas the deities of the oldest centres of culture Eridu, Uruk and Nippur always remained the most highly venerated gods Ea, Ami and Enlil (mostly called Bel, i.e. 'Lord'), after 2000 BG the political hegemony of Babylon made its local god Marduk the senior god of the pantheon. When in later centuries Assyria acquired the overlordship of the Mesopotamian world, Ashur took this place. When combined into a hierarchy, the gods lost some of their local character, and their character as personifications of natural powers became more prominent. The deities Sin of Ur and Shamash of Sippara in later times were always venerated as the moon god and the sun god.

এমন সাম্রাজ্যে রাজা ছিলেন আইনপ্রণেতা, প্রধান বিচারক এবং কর্মকর্তাদের নিয়ে গঠিত সরকারের প্রধান। বেসামরিক বিভাগগুলোর নেতা হিসেবে এই কর্মকর্তারা সামরিক শ্রেণির পাশাপাশি দ্বিতীয় একটি শাসকশ্রেণি গঠন করতেন। সাধারণত এই শ্রেণি গঠিত হতো পুরোহিতদের নিয়ে। তাঁরা আগে স্থানীয় বুদ্ধিবৃত্তিক নেতা ছিলেন, এখন একটি দাপ্তরিক স্তরবিন্যাসে সংগঠিত হলেন। রাষ্ট্র ও সমাজের আধ্যাত্মিক নেতৃত্ব ছিল পুরোহিততন্ত্রের হাতে। উৎপাদনপ্রক্রিয়ার জন্য প্রয়োজনীয় তাত্ত্বিক ও সাধারণ জ্ঞান তাঁদের কাছেই ছিল। এটাই ছিল তাঁদের মর্যাদা ও সামাজিক ক্ষমতার উৎস। যেখানে কৃষিই ছিল প্রধান পেশা, সেখানে পঞ্জিকা ও ঋতুর জ্ঞান ছিল তাঁদের এখতিয়ারে।

সে যুগে ধর্মও রাষ্ট্র ও সমাজের সঙ্গে একাকার ছিল, এবং ধর্মও কেন্দ্রীভূত হলো। প্রধান নগরগুলোর স্থানীয় দেবতারা একজন প্রধান দেবতার অধীনে একটি দেবমণ্ডলীতে একত্রিত হলেন। অন্যদিকে ছোট বা বিজিত স্থানগুলোর স্থানীয় দেবতারা নিম্নস্তরের আত্মায় নেমে গেলেন, কিংবা অন্য দেবতাদের মধ্যে মিশে গেলেন। যেমন প্রাচীন ব্যাবিলোনিয়ায় দেবী ইন্নিনা, নিসাবা ও নানা পরবর্তীকালের ইশতারের সঙ্গে এক হয়ে যান। দ্রুত বেড়ে ওঠা ব্যাবিলন নগরীর শহরতলিতে পরিণত হয় বোরসিপ্পা। এর ধর্মতাত্ত্বিক প্রকাশ ঘটে এভাবে যে বোরসিপ্পার নগরদেবতা নাবো মারদুকের পুত্র হয়ে যান। সংস্কৃতির প্রাচীনতম কেন্দ্র এরিদু, উরুক ও নিপ্পুরের দেবতারা, অর্থাৎ এয়া, আনু ও এনলিল (যাঁকে প্রায়ই বেল, অর্থাৎ "প্রভু" বলা হতো), সবসময়ই সর্বাধিক পূজিত দেবতা ছিলেন। তবু খ্রিস্টপূর্ব ২০০০ সালের পর ব্যাবিলনের রাজনৈতিক আধিপত্যের ফলে তার স্থানীয় দেবতা মারদুক দেবমণ্ডলীর প্রধান দেবতা হয়ে ওঠেন। পরের শতাব্দীগুলোতে আসিরিয়া যখন মেসোপটেমীয় জগতের ওপর কর্তৃত্ব লাভ করে, তখন আশুর এই স্থান নেন। স্তরবিন্যাসে একত্রিত হওয়ার পর দেবতারা তাঁদের স্থানীয় চরিত্র কিছুটা হারান, আর প্রাকৃতিক শক্তির মূর্ত রূপ হিসেবে তাঁদের চরিত্র বেশি প্রকট হয়ে ওঠে। যেমন উরের দেবতা সিন এবং সিপ্পারার দেবতা শামাশ পরবর্তীকালে সবসময় চন্দ্রদেব ও সূর্যদেব হিসেবে পূজিত হয়েছেন।

এমন একটি গোষ্ঠীর উত্থান ঘটল, যাদের শাসকশ্রেণি হিসেবে আর কঠোর পরিশ্রম করে জীবন নিরাপদ রাখতে হতো না। এর ফলে অস্তিত্বের নতুন পরিস্থিতি তৈরি হলো। সমাজকাঠামো আরও জটিল হলো, আর আধ্যাত্মিক নেতৃত্বের জন্য উচ্চতর যোগ্যতার প্রয়োজন হলো এবং তার দাবিও বিস্তৃত হলো। ব্যবসা-বাণিজ্য নতুন বস্তুগত ও আত্মিক চাহিদার জন্ম দিল। রাজা ও অভিজাতদের সম্পদ ও বিলাসিতার সঙ্গে এল শিল্প ও বিজ্ঞানের প্রতি অনুরাগ। এভাবে পৃথিবীতে প্রথমবারের মতো নতুন সমাজকাঠামোর পাশাপাশি সংস্কৃতির উচ্চতর রূপের উদ্ভব হলো, যা প্রাগৈতিহাসিক বর্বরতার সবচেয়ে উন্নত রূপকেও ছাড়িয়ে গেল। শুরু হলো সভ্যতার যুগ।

এর সঙ্গে জ্যোতির্বিজ্ঞানেরও উচ্চতর বিকাশ ঘটল। এর উৎস ছিল সরাসরি কালগণনার প্রয়োজন, বিশেষ করে চান্দ্র পঞ্জিকাকে সৌর বছরের সঙ্গে খাপ খাওয়ানোর সমস্যা।

চাঁদের একটি পর্যায়কাল গড়ে ২৯.৫৩০৫৯ দিন। একটি সৌর বছর ৩৬৫.২৪২২০ দিন, যা ১২টি চান্দ্র পর্যায়কালের (মোট ৩৫৪.৩৬৭১ দিন) চেয়ে ১১ দিন বেশি। তিন বছর পর চান্দ্র পঞ্জিকা সূর্যের অগ্রগতি থেকে ৩৩ দিন পিছিয়ে পড়ে। সূর্যের সঙ্গে সামঞ্জস্য রাখতে প্রতি তৃতীয় বছরে, কখনো আরও ঘন ঘন, একটি বাড়তি মাস যোগ করতে হতো, যাতে বছরে ১২টির বদলে ১৩টি মাস হয়। পঞ্জিকার সমস্যাটি হলো এমন একটি বড় পর্যায়কাল খুঁজে বের করা, যা মাস ও বছর উভয়েরই সাধারণ গুণিতক। তাহলে এই পর্যায়কাল শেষে সূর্য ও চাঁদ আবার একই পারস্পরিক অবস্থানে ফিরে আসবে। অবশ্য নিখুঁত কোনো সাধারণ গুণিতক নেই, তবে কমবেশি সন্তোষজনক আসন্ন মান পাওয়া যায়।

সৌর ও চান্দ্র পর্যায়কাল সম্পর্কে আমাদের নির্ভুল আধুনিক জ্ঞান দিয়ে এমন আসন্ন পঞ্জিকা-পর্যায়কাল তাত্ত্বিকভাবে বের করা যায়। এর জন্য দুটির অনুপাতকে একটি অবিচ্ছিন্ন ভগ্নাংশে (continued fraction) রূপান্তর করে তার ক্রমিক আসন্ন মানগুলো লিখতে হয়। এভাবে আমরা ৩৭/৩ ও ৯৯/৮ পাই। প্রথমটি দেখায় যে ৮ বছর প্রায় ৯৯ মাসের সমান (যথাক্রমে ২৯২১.৯৪ ও ২৯২৩.৫৩ দিন)। তাই এই ৮ বছরের মধ্যে ৩টিতে ১৩তম মাস থাকতে হবে, আর ৫টিতে থাকবে ১২টি মাস। তবু এই আসন্ন মান খুব ভালো নয়। মাত্র ২৪ বছর পরই চান্দ্র তারিখ সৌর ঋতুর তুলনায় ৫ দিন পিছিয়ে পড়বে। এর চেয়ে অনেক বেশি নির্ভুল হলো ১৯ বছরের একটি পর্যায়কাল, যাতে ২৩৫টি মাস থাকে (যথাক্রমে ৬৯৩৯.৬০ ও ৬৯৩৯.৬৯ দিন)। এখানে ১৩তম মাস সাতবার যোগ করতে হয়। অবশ্যই আদিতম যুগে মানুষের পর্যায়কাল সম্পর্কে এই নিখুঁত জ্ঞান ছিল না। তাদের জন্য একটি ভালো পঞ্জিকা-পর্যায়কাল খুঁজে পাওয়া ছিল কঠিন এক ব্যবহারিক সমস্যা। সৌর ও চান্দ্র গণনাকে খাপ খাওয়ানোর শ্রমসাধ্য প্রক্রিয়ার মাধ্যমেই কেবল এর সমাধান সম্ভব ছিল। তাই এই সমস্যা আকাশের ঘটনাবলিকে আরও যত্নের সঙ্গে পর্যবেক্ষণ করার চালিকাশক্তি হিসেবে কাজ করেছিল।

প্রশ্ন উঠতে পারে, এত নির্ভুলতার প্রয়োজন কেন ছিল? এটা তো কৃষির প্রয়োজনের চেয়ে অনেক বেশি, আর আবহাওয়ার ওঠানামার কারণে কৃষিকাজ এমনিতেই বেশ অনিয়মিত। তবে আমাদের মনে রাখতে হবে, সে যুগে কৃষিকাজের সঙ্গে ধর্মীয় আচার ও উৎসব জড়িত থাকত। অন্য সব বড় ও গুরুত্বপূর্ণ সামাজিক ঘটনার মতো কৃষি-উৎসবগুলোও একই সঙ্গে ছিল ধর্মীয় অনুষ্ঠান। প্রধান প্রাকৃতিক ও সামাজিক শক্তিগুলোর প্রতিনিধি হিসেবে দেবতারা মানুষের জীবনে অংশ নিতেন। সামাজিকভাবে যা প্রয়োজনীয় বা উপযুক্ত ছিল, তা দেবতাদের আদেশে পরিণত হতো এবং আচার-অনুষ্ঠানে কঠোরভাবে নির্দিষ্ট করে দেওয়া হতো। প্রকৃতিগতভাবে যা একটি নির্দিষ্ট ঋতুতে ঘটত, যেমন ফসল ঘরে তোলার উৎসব, ধর্মীয় অনুষ্ঠান হিসেবে তা একটি নির্দিষ্ট তারিখে বেঁধে দেওয়া হতো, যেমন চাঁদের কোনো একটি কলায়। দেবতাদের সেবায় কোনো অবহেলার সুযোগ ছিল না, সেখানে আচার-অনুষ্ঠান নিখুঁতভাবে পালন করতে হতো। পঞ্জিকা মূলত ছিল আচার-অনুষ্ঠানের কালানুক্রমিক বিন্যাস। এই কারণে পঞ্জিকা দায়িত্বপ্রাপ্ত কর্মকর্তাদের, বিশেষ করে পুরোহিতদের, নিরন্তর যত্নের বিষয় হয়ে উঠেছিল। আবার শুভ সময়ের জ্ঞানে তাঁদের একচেটিয়া অধিকার থাকায় পঞ্জিকা একই সঙ্গে তাঁদের সামাজিক ক্ষমতারও উৎস ছিল।

এ ধরনের ধর্মীয়-কৃষিভিত্তিক চর্চা থেকে কীভাবে চান্দ্র-সৌর পঞ্জিকা প্রতিষ্ঠিত হলো, তার কয়েকটি চমকপ্রদ দৃষ্টান্ত আমাদের কাছে এসে পৌঁছেছে। ইহুদি পঞ্জিকার উৎপত্তি মরুভূমিতে, তাই এটি চাঁদের ওপর ভিত্তি করে তৈরি হয়েছিল। কিন্তু ইহুদিরা কেনান দেশে পৌঁছানোর পর কৃষিই প্রধান পেশা হয়ে উঠল, আর পঞ্জিকাকে সূর্যের সঙ্গে খাপ খাওয়াতে হলো। বসন্তকালে মাসোথ ছিল ফসল তোলার উৎসব। এতে যবের প্রথম আঁটিগুলো নিবেদন করা হতো, আর প্রথম শস্য দিয়ে খামিরবিহীন রুটি তৈরি করা হতো। এই উৎসব যাযাবরদের পেসাখ উৎসবের সঙ্গে মিশে যায়। পেসাখে বসন্তের পূর্ণিমায় সদ্যোজাত মেষশাবক জিহোভার উদ্দেশে উৎসর্গ করা হতো। এভাবে উৎসবটি প্রথম মাস নিসানের পূর্ণিমায় স্থির হলো।

কীভাবে এটি করা হতো, তা গিনৎসেলের কালগণনা-বিষয়ক বিখ্যাত পাঠ্যবইয়ে পড়া যায়: "শেষ মাসের শেষ দিকে পুরোহিতেরা মাঠে গিয়ে ফসলের অবস্থা পরিদর্শন করতেন, পরের দুই সপ্তাহে যব পাকবে বলে আশা করা যায় কি না তা দেখতে। যদি দেখতেন যে পাকবে, তাহলে পরবর্তী অমাবস্যা থেকে শুরু হওয়া মাসেই মাসোথ-পেসাখ নির্ধারণ করা হতো। কিন্তু যদি পাকার আশা না থাকত, তাহলে উৎসবের মাস এক চান্দ্র পর্যায়কাল পিছিয়ে দেওয়া হতো। এর ভিত্তিতেই অন্যান্য উৎসবের দিন নির্ধারিত হতো।"

ইসরায়েলিরা যতদিন প্যালেস্টাইনে একসঙ্গে বাস করত, ততদিন এই অভিজ্ঞতাভিত্তিক পদ্ধতি তাদের কাজে লেগেছিল। পরে তারা ছড়িয়ে পড়লে এটি আর ব্যবহার করা গেল না। তবে ততদিনে জ্যোতির্বিজ্ঞান এতটা এগিয়ে গিয়েছিল যে তারা প্রতিবেশীদের কাছ থেকে পঞ্জিকা-পর্যায়কালের জ্ঞান গ্রহণ করতে পারল। তখন ১৯ বছরের পর্যায়কাল ইহুদি পঞ্জিকার ভিত্তি হয়ে উঠল। এতে ১২ মাসের ১২টি বছর আর ১৩ মাসের ৭টি বছর একটি নির্দিষ্ট ক্রমে পালাক্রমে আসে।

আরেকটি দৃষ্টান্ত হলো মুহাম্মদের পূর্ববর্তী আরবের ঘটনাপ্রবাহ। পবিত্র মাসগুলোতে রক্তের প্রতিশোধ স্থগিত রাখা হতো এবং কাফেলাগুলো নিরাপদে চলাচল করতে পারত, কারণ অর্থনৈতিক জীবনের জন্য এটা দরকার ছিল। গোত্রগুলোর মধ্যে অন্তহীন বিবাদ ছিল খুবই স্বাভাবিক ব্যাপার। গোত্রের সদস্যদের মধ্যে সংহতির একটি আদিম বিচারিক রূপ হিসেবে রক্তের প্রতিশোধও প্রয়োজনীয় ছিল। কিন্তু এর সীমাহীন দাপট চললে বেচাকেনা আর খাদ্যের জোগান অসম্ভব হয়ে পড়ত। মাসউদী বলেছেন: "সফর মাসের নাম এসেছে ইয়েমেনের বাজার থেকে… এখানে আরবরা তাদের শস্য কিনত, আর যে তা করতে ভুলত সে না খেয়ে মারা যেত।"

পবিত্র মাস ছিল বাজারের মাস। চারদিক থেকে কাফেলাগুলো বড় বড় বাজার-শহরে যেত, বিশেষ করে মক্কায়। মক্কা তার অতুলনীয় শীতল কূপ জমজমের জন্য বিখ্যাত ছিল, যা আরবে তীর্থকেন্দ্র হিসেবে মক্কার গুরুত্বের উৎস। এখানে গোত্রগুলো মিলিত হতো ও আলাপ-আলোচনা করত, আর এই মিলনস্থলই রাজনৈতিক ও ধর্মীয় কেন্দ্র এবং রাজধানীতে পরিণত হলো। তাই পবিত্র মাসটি সবসময় বছরের একই সময়ে পড়া জরুরি ছিল, যখন ফসল তোলা হয়ে গেছে এবং পণ্য হাতে আছে।

পরবর্তীকালের মুসলিম পণ্ডিত আল-বিরুনী লিখেছেন: "জাহেলি যুগে আরবরা তাদের মাসগুলোকে সেভাবেই গণনা করত যেভাবে মুসলমানরা এখন করে, আর তাদের হজ বছরের চারটি ঋতু জুড়ে ঘুরে বেড়াত। তবে পরে তারা স্থির করল যে হজ এমন সময়ে হবে যখন তাদের পণ্য, চামড়া ও ফল বাজারে বিক্রির জন্য প্রস্তুত থাকে। তাই তারা একে অচল করতে চাইল, যাতে এটি সবসময় সবচেয়ে প্রাচুর্যের ঋতুতে পড়ে। এজন্য হিজরতের ২০০ বছর আগে তারা ইহুদিদের কাছ থেকে অধিমাস যোগ করার পদ্ধতি শিখল। ইহুদিদের মতো করেই তারা এটি প্রয়োগ করল: তাদের বছর আর সৌর বছরের পার্থক্য যখন বেড়ে এক মাসের সমান হতো, তখন তা বছরের মাসগুলোর সঙ্গে যোগ করত। হজের অনুষ্ঠান শেষে কালাম্মাসরা (এই দায়িত্বপ্রাপ্ত একটি গোত্রের শেখরা) সামনে এসে জনতার উদ্দেশে বক্তৃতা দিতেন। তাঁরা পরের মাসটিকে চলতি মাসের নাম দিয়ে একটি মাস যোগ করতেন। জনতা হাততালি দিয়ে কালাম্মাসদের সিদ্ধান্তে সম্মতি জানাত। এই প্রক্রিয়াকে তারা বলত নাসি, অর্থাৎ স্থানান্তর, কারণ প্রতি দ্বিতীয় বা তৃতীয় বছরে বছরের শুরু সরে যেত… মানাজিলের উদয় ও অস্ত দেখে তারা এটি বিচার করতে পারত। মক্কা থেকে মদিনায় নবীর হিজরত পর্যন্ত এই ব্যবস্থাই চালু ছিল।"

হিজরতের পর নবম বছরে মুহাম্মদ এই স্থানান্তর নিষিদ্ধ করেন। হয়তো কালাম্মাসদের দায়িত্ব কেড়ে নিয়ে তাঁদের আধ্যাত্মিক ক্ষমতা ভেঙে দেওয়া এবং ইহুদিদের থেকে নিজেকে আরও স্পষ্টভাবে আলাদা করাই ছিল এর উদ্দেশ্য। তাই এখন মুসলিম পঞ্জিকা ৩৫৪ দিনের চান্দ্র বছরের ওপর ভিত্তি করে চলে। প্রতিটি বছরে ১২টি চান্দ্র মাস থাকে, আর ৩৩ বছরে এগুলো সব ঋতু ঘুরে আসে। এখানে আমরা এমন একটি পঞ্জিকা পাই যা আনুষ্ঠানিক বিধানের ফলে সামাজিক চর্চা থেকে বিচ্ছিন্ন হয়ে গেছে। এটি আদিম বেদুইন জীবনধারার এক প্রস্তরীভূত ধর্মীয় ঐতিহ্যে পরিণত হয়েছে।

এটি দেখাতে পারে যে ভিন্ন ভিন্ন অর্থনৈতিক পরিস্থিতি ও রাজনৈতিক ইতিহাসের ফলে বিভিন্ন জাতির মধ্যে পঞ্জিকা ও জ্যোতির্বিজ্ঞান কীভাবে ভিন্ন ভিন্ন পথে বিকশিত হয়েছে।

‌% ch3.tex


\chapter{প্রাচীন ব্যাবিলনীয় নক্ষত্রবিদ্যা}
\epigraph{``\textit{
খেলিছ এ বিশ্ব লয়ে \\
বিরাট শিশু আনমনে। \\
প্রলয় সৃষ্টি তব পুতুল খেলা \\
নিরজনে প্রভু নিরজনে।। \\
}''}{\textit{নজরুল ইসলাম}}

যে প্রাচীনতম কাল থেকে তথ্য আমাদের কাছে পৌঁছেছে, তখন ইউফ্রেটিস ও টাইগ্রিসের মধ্যবর্তী সমভূমি নিয়ে গঠিত শিনার দেশে দুটি স্বতন্ত্র জাতির বাস ছিল। উত্তরাংশ আক্কাদে ছিল একটি সেমিটিক জাতি, আর দক্ষিণাংশে সুমেরীয়রা। ভাষা, চেহারা ও রীতিনীতিতে এই দুই জাতি একেবারেই আলাদা ছিল। সুমেরীয়রা ছিল আদি অধিবাসী বা প্রথম অভিবাসী। অন্যদিকে সেমিটিরা পরে পশ্চিমের মরুভূমি থেকে ক্রমবর্ধমান সংখ্যায় এখানে আসে। সুমেরীয় ভাষার সঙ্গে সেমিটিক বা ইন্দো-ইউরোপীয় ভাষার কোনো সম্পর্ক দেখা যায় না, বরং তুরানীয় ভাষার সঙ্গে সাদৃশ্য আছে। যা-ই হোক, এই প্রাচীনতম সংস্কৃতির জাতির নৃতাত্ত্বিক পরিচয় অজানা।

খ্রিস্টপূর্ব ৩০০০ সালের পরের শতাব্দীগুলোতে দক্ষিণের সুমেরীয় নগরগুলোতে সংস্কৃতি সবচেয়ে উন্নত ছিল: সমুদ্রতীরের এরিদু, মরুভূমির প্রান্তে উর, এবং উরুক, লাগাশ, নিপ্পুর ও লারসা। সাধারণত এগুলোর একটির হাতে আধিপত্য থাকত। উত্তরের নগরগুলোর (আগাদে, সিপ্পার, বোরসিপ্পা, ব্যাবিলন) সেমিটিরা এই সংস্কৃতি গ্রহণ করে, এবং আরও অভিবাসনের পর ক্রমশ প্রধান জাতিতে পরিণত হয়। খ্রিস্টপূর্ব ২৫০০ সালের দিকে আগাদের সারগন এবং পরে তাঁর পুত্র নারাম-সিন যখন সমগ্র মেসোপটেমিয়া শাসন করতেন, তখন সামরিক অংশ ছিল সেমিটিদের নিয়ে গঠিত, আর লিপিকার ও বেসামরিক কর্মকর্তারা ছিলেন সুমেরীয়। পরের শতাব্দীগুলোতে কর্তৃত্ব আবার দক্ষিণে ফিরে যায়। সেখানে লাগাশ ও উরের শাসকেরা নিজেদের "সুমের ও আক্কাদের রাজা" বলতেন। খ্রিস্টপূর্ব ২০০০ সালের পর পশ্চিম থেকে আমোরীয় গোত্রগুলোর অভিবাসনের ফলে সেমিটিরা নিশ্চিত প্রাধান্য লাভ করে। তখন একটি রাজবংশের অধীনে, যার সবচেয়ে বিখ্যাত রাজা হাম্মুরাবি, ব্যাবিলন নগরী একটি বিশাল সাম্রাজ্যের রাজধানী এবং বাণিজ্য ও সংস্কৃতির কেন্দ্র হয়ে ওঠে।

সুমেরীয়রাই কীলকাকার লিপির উদ্ভাবক। এই লিপিতে একটি স্বরবর্ণ এবং এক বা দুটি ব্যঞ্জনবর্ণ নিয়ে গঠিত প্রতিটি ধ্বনির জন্য একটি বিশেষ চিহ্ন ছিল। নরম কাদার ফলকে একটি লেখনী, অর্থাৎ ধাতব কীলক, দিয়ে ছাপ দিয়ে চিহ্নগুলো তৈরি করা হতো। ছাপগুলো এক পাশে চওড়া, অন্য পাশে সরু হতো। তারপর ফলকটি আগুনে পুড়িয়ে শক্ত করা হতো। চিহ্ন দিয়ে অক্ষর প্রকাশের এই পদ্ধতি সুমেরীয়ের মতো একটি সংযোজনশীল (agglutinating) ভাষার জন্য পুরোপুরি উপযোগী ছিল, যেখানে আলাদা আলাদা শব্দ ও ধাতু কেবল পরপর জুড়ে দেওয়া হতো। সেমিটিরা এই কীলকাকার লিপি গ্রহণ করে, যদিও ধাতুর রূপান্তরভিত্তিক তাদের ভাষার জন্য এটি বেশ অসুবিধাজনক ছিল। ব্যাবিলন যখন সমগ্র নিকট প্রাচ্যের সাংস্কৃতিক কেন্দ্র হয়ে উঠল, তখন এর ভাষা ও কীলকাকার লিপি মিশর ও এশিয়া মাইনর পর্যন্ত আন্তর্জাতিক যোগাযোগে ব্যবহৃত হতো।

কীলকাকার লিপি এবং তাতে লেখা ভাষাগুলোর পাঠোদ্ধার ছিল উনিশ শতকের অন্যতম বড় কৃতিত্ব। এর ফলে ইতিহাস ও সংস্কৃতির একটি সম্পূর্ণ হারিয়ে যাওয়া যুগ আলোয় এল। আগের জ্ঞান বেশিরভাগই এসেছিল গ্রিক উৎস থেকে, যা ছিল খণ্ডিত ও উপাখ্যাননির্ভর, এবং তা খ্রিস্টপূর্ব ৭০০ সালের আগে যেত না। ১৮৪৬ সালে ফরাসি কনসাল বোত্তার প্রথম খননকাজে অনুপ্রাণিত হয়ে হেনরি লেয়ার্ড প্রাচীন নিনেভের স্থানে খনন শুরু করেন। তিনি অপূর্ব সব ভাস্কর্য খুঁড়ে বের করেন: শিকারের দৃশ্যের উদ্‌গত খোদাই, ডানাওয়ালা ষাঁড় ও ড্রাগন। চারপাশে ছড়িয়ে থাকা কীলকাকার লিপিযুক্ত অসংখ্য ভাঙা টুকরো দেখে তিনি বিস্মিত হন। সেগুলোর মূল্য আন্দাজ করে তিনি শিল্পকর্মগুলোর সঙ্গে বাক্সভর্তি এই টুকরোও ব্রিটিশ মিউজিয়ামে পাঠান। বহু বছর পর জর্জ স্মিথ যখন মহাপ্লাবনের কাহিনিসংবলিত কিছু লেখার পাঠোদ্ধার করেন, তখন এগুলোর গুরুত্ব বোঝা যায়। তারপর যত বেশি সম্ভব টুকরো সংগ্রহের জন্য বিশেষ অভিযান পাঠানো হয়। দেখা গেল, এগুলো রাজা আশুরবানিপালের "গ্রন্থাগারের" অবশেষ। এতে নতুন শিলালিপির পাশাপাশি আগের শতাব্দীগুলোর লেখার অনুলিপি, অভিধান ও অভিধান-সংক্রান্ত উপকরণ ছিল। এগুলো প্রাচীন মেসোপটেমীয় সংস্কৃতির চর্চায় (যাকে তখন থেকে "অ্যাসিরিওলজি" বলা হয়) জোরালো প্রেরণা জোগায়।

শুরুতে দুটি সম্পূর্ণ ভিন্ন ভাষার মিশ্রণ অ্যাসিরিওলজিস্টদের বিভ্রান্ত করেছিল। সাধারণত প্রতিটি চিহ্ন ধ্বনি হিসেবে ব্যবহৃত হতো, অর্থাৎ সেমিটিক ভাষার জন্য সুমেরীয় থেকে একটি অক্ষর হিসেবে ধার নেওয়া হতো। কিন্তু প্রায়ই চিহ্নটি সুমেরীয় ভাষায় তার অর্থ অনুযায়ী তথাকথিত "ভাবচিত্র" (ideogram) হিসেবে, অর্থাৎ কোনো বস্তু বা ধারণার প্রতীক হিসেবে ব্যবহৃত হতো। এভাবে দুটি ভাষারই রহস্য ধীরে ধীরে একসঙ্গে উন্মোচিত হলো। অন্যান্য প্রাচীন নগরের ধ্বংসাবশেষ থেকে খুঁড়ে পাওয়া লেখা এই সাহিত্যিক খণ্ডগুলোকে পূর্ণতা দিল, আর এগুলো প্রাচীনকালের ইতিহাস ও সংস্কৃতিকে ক্রমশ আরও স্পষ্ট করে তুলল।

এর মাধ্যমে জ্যোতির্বিজ্ঞানের একটি অজানা প্রাচীন ইতিহাসও আলোয় এল। এই পুনরাবিষ্কৃত প্রাচীন জগতে জ্যোতিষ্কগুলো অন্য যেকোনো দেশ বা যুগের চেয়ে বেশি গুরুত্ব পেয়েছিল। ছোট ছোট বিক্ষিপ্ত খণ্ড থেকে এই প্রাচীন জ্ঞান পুনর্গঠন করা অবশ্যই কঠিন ছিল। অস্পষ্ট অনুমান, ঝুঁকিপূর্ণ ধারণা, ভুল ব্যাখ্যা আর অসমর্থনযোগ্য তত্ত্বের মধ্য দিয়েই কেবল ধীরে ধীরে সত্যের কাছাকাছি পৌঁছানো গেছে। বহু বছর ধরে হুগো উইংক্লারের তত্ত্ব, তথাকথিত "প্যান-ব্যাবিলনবাদ", খুব খ্যাতি পেয়েছিল। এই তত্ত্ব ঘোষণা করত যে প্রাচীনতম কালে, খ্রিস্টপূর্ব ৩০০০ থেকে ২০০০ সালের মধ্যে, সেখানে ইতিমধ্যেই একটি অত্যন্ত উন্নত জ্যোতির্বিজ্ঞান ছিল। এর ভিত্তি ছিল আকাশের পর্যায়কালগুলোর পূর্ণাঙ্গ জ্ঞান, এবং বিষুবের অয়নচলনের (precession of the equinoxes) ফলে তারামণ্ডলের চেহারা কীভাবে সরে যায় তার জ্ঞান। তত্ত্বটি দাবি করত যে এই আদিম বিশ্বধারণা পার্থিব ও আকাশের ঘটনাবলির মধ্যে ঘনিষ্ঠ সম্পর্ক ঘোষণা করত বলে এর চরিত্র ছিল "নাক্ষত্রিক পুরাণের"। এটিই পরবর্তী সব প্রাচ্য ও গ্রিক চিন্তাধারার উৎস, এবং আধুনিক কাল পর্যন্ত ইউরোপের দূরবর্তী অঞ্চলেও কিংবদন্তি, লোকজ্ঞান ও রীতিনীতি নির্ধারণ করেছে। পরে লেখাগুলো যত্নের সঙ্গে অধ্যয়ন করা হলে এই আকর্ষণীয় কিন্তু অত্যন্ত কল্পনাপ্রসূত তত্ত্ব টেকেনি। ব্যাবিলনীয় জ্যোতির্বিজ্ঞান সম্পর্কে আমাদের বর্তমান প্রকৃত জ্ঞান মূলত তিনজন জেসুইট পণ্ডিতের কাজের ওপর দাঁড়িয়ে আছে: অ্যাসিরিওলজিস্ট জে. এন. স্ট্রাসমায়ার এবং জ্যোতির্বিদ জে. এপিং ও এফ. এক্স. কুগলার।

প্রথম ব্যাবিলনীয় রাজ্যের যুগ ছিল অর্থনৈতিক ও রাজনৈতিক শক্তি এবং সাংস্কৃতিক জীবনের চূড়ান্ত বিকাশের সময়। বাণিজ্য ও শিল্প সমৃদ্ধ হয়েছিল। অত্যন্ত গুরুত্বপূর্ণ বাণিজ্যনগরী ব্যাবিলন কেবল একটি বিশাল মেসোপটেমীয় সাম্রাজ্যের রাজধানীই ছিল না, বরং সমগ্র নিকট-এশিয়ায় আলো ছড়ানো একটি সাংস্কৃতিক কেন্দ্রও ছিল। এখানে পূর্ববর্তী শতাব্দীগুলোর সুমেরীয় সংস্কৃতির ফল পূর্ণতা পায়, আর ধর্মতত্ত্বের ব্যবস্থা এমন রূপ নেয় যা পরবর্তীকালেও বহাল থাকে। বিখ্যাত "হাম্মুরাবির আইনসংহিতা" থেকে আমরা নাগরিক জীবনের কাঠামো ভালোভাবে জানি। এটি একটি পাথরে খোদাই করা, যা উনিশ শতকে সুসায় খনন করে পাওয়া যায়। খুঁড়ে পাওয়া অসংখ্য মাটির ফলকের মধ্যে জমি কেনাবেচা, ঋণ দেওয়া, ভাড়া ও চাকরি সংক্রান্ত বহু দেওয়ানি চুক্তি পাওয়া গেছে। এগুলো মন্দিরে জমা রেখে সংরক্ষণ করা হতো, যেন মন্দিরগুলো নোটারি অফিস। এগুলোর মাধ্যমে আমরা পঞ্জিকা সম্পর্কে এবং রাজাদের নাম ও তারিখের ক্রম সম্পর্কে নিশ্চিত হতে পেরেছি, যা ইতিহাসের কাঠামো হিসেবে কাজ করে।

উরের রাজা দুংগির আমলে প্রচলিত বারো মাসের সুমেরীয় নামগুলো বিভিন্ন জায়গায় বিভিন্ন ছিল। এগুলোর মধ্যে দেখা যায়, চতুর্থ মাসের নাম গঠিত হয়েছে "বীজ" ও "হাত"-এর চিহ্ন দিয়ে, একাদশ মাসের নাম "শস্য" ও "কাটা"-র চিহ্ন দিয়ে, আর দ্বাদশ মাসের নাম "শস্য" ও "ঘর"-এর চিহ্ন দিয়ে। এভাবে ঋতুটি স্পষ্টভাবে নির্দেশিত হয়েছে। তাই আমরা বুঝতে পারি কেন সেখানে অধিমাস যোগ করার জন্য একাদশ মাসটি দ্বিগুণ করা হতো। ব্যাবিলনের উত্থানের পর সেমিটিক নামগুলো চালু হয়: নিসান্নু, আইয়ারু, সিমান্নু, দুজু, আবু, উলুলু, তিশরিতু, আরাখ-সাম্না, কিসলিমু, তেবিতু, শাবাতু, আদারু। জ্যোতির্বৈজ্ঞানিক লেখায় এগুলোকে একটি করে চিহ্ন দিয়ে নির্দেশ করা হয়, যা তাদের আগের সুমেরীয় নামের প্রথম অক্ষর।

প্রাচীন ব্যাবিলনীয় রাজ্যে অধিমাসটি ছিল বছরের শেষে একটি দ্বিতীয় আদারু। দেওয়ানি চুক্তির তারিখগুলোতে এই অধিমাস যোগের অনিয়ম দেখা যায়। এ থেকে বোঝা যায় যে এগুলো অভিজ্ঞতার ভিত্তিতে করা হতো, ফসল পাকার অবস্থা অনুযায়ী বা অন্য কোনো কারণে প্রয়োজন মনে হলে। পঞ্জিকা অনেক বেশি পিছিয়ে পড়লে এমনকি পরপর দুই বছরও ১৩ মাসের হয়েছে। কয়েকবার, স্পষ্টতই কোনো বিচ্যুতি দ্রুত সংশোধনের প্রয়োজন হলে, ষষ্ঠ মাসটি দ্বিগুণ করা হয়েছিল। এমন একটি ঘটনার দলিল আছে, যাতে লেখা: "হাম্মুরাবি এভাবে বলছেন: 'যেহেতু বছরটি ভালো নয়, তাই পরের মাসটিকে দ্বিতীয় উলুলু হিসেবে লিখতে হবে। তিশরিতুর ২৫ তারিখে ব্যাবিলনে দশমাংশ পৌঁছে দেওয়ার বদলে দ্বিতীয় উলুলুর ২৫ তারিখে পৌঁছে দেওয়ার ব্যবস্থা করো।'" রাজদরবারের জন্য খাদ্য সরবরাহ অবশ্যই এক মাস পিছিয়ে দেওয়া যেত না।

পঞ্জিকার জন্য কিছু আকাশের ঘটনা পর্যবেক্ষণ করা জরুরি ছিল। মাসের প্রথম দিনটি সঠিকভাবে নির্ধারণ করতে নতুন চাঁদকে তার প্রথম আবির্ভাবেই ধরতে হতো। এটি খুব কঠিন ছিল না। সেই অপূর্ব উজ্জ্বল জলবায়ুতে (শীতের কয়েকটি মাস আর মাঝে মাঝে বালুঝড় বাদে) বিস্তীর্ণ সমভূমির এক দিগন্ত থেকে আরেক দিগন্ত পর্যন্ত স্পষ্ট দেখা যেত। পুরোহিত-জ্যোতির্বিদেরা তাঁদের ধাপযুক্ত মিনার থেকে সান্ধ্য আকাশে চাঁদের প্রথম সরু ফালি সহজেই দেখতে পেতেন। আনুষ্ঠানিক প্রয়োজনে তাঁদের পূর্ণিমার দিকেও নজর রাখতে হতো। আর পূর্বাভাসের জন্য মাসের শেষ দিকে প্রভাতী বাঁকা চাঁদের শেষ আবির্ভাবও পর্যবেক্ষণ করতে হতো। এসব করতে গিয়ে তাঁরা কি নীরবে নিজেদের পথ আঁকতে থাকা আর প্রতি মাসে এগিয়ে যাওয়া নক্ষত্রগুলো লক্ষ করেননি, এমন কি কল্পনা করা যায়? কিংবা তাদের মাঝে মাঝে মাঝে দেখা দেওয়া উজ্জ্বল গ্রহগুলো, বিশেষ করে অতুলনীয় সন্ধ্যাতারা, কি তাঁদের মুগ্ধ করেনি? এগুলো নিয়ে লেখা খুব কম। অত পরিষ্কার আকাশ থাকা সত্ত্বেও আদি ব্যাবিলনীয়দের মন নক্ষত্রে ততটা নিবিষ্ট ছিল না, যতটা প্রায়ই ধারণা করা হয়েছে। তবে খুবই সম্ভব, এবং কিছু লেখাও এটি সমর্থন করে, যে চাঁদের এই নিয়মিত পর্যবেক্ষণ ধীরে ধীরে নক্ষত্র নিয়ে ব্যবহারিক আগ্রহ বাড়িয়ে তুলেছিল।

নতুন চাঁদ পর্যবেক্ষণে প্রথম যে ঘটনাটি ধরা পড়ে, তা হলো বছরের পরিক্রমায় তারামণ্ডলগুলোর নিয়মিত অগ্রসর হওয়া ও সরে যাওয়া। পশ্চিমের সান্ধ্য আকাশে দৃশ্যমান নক্ষত্রগুলো ঋতুর বৈশিষ্ট্য বহন করে। তাই এগুলো দিয়ে পঞ্জিকা ও অধিমাস যাচাই করা যায়। পূর্বের প্রভাতী আকাশে নক্ষত্রের প্রথম আবির্ভাবের ক্ষেত্রেও একই কথা খাটে। তাই নক্ষত্রের নিয়মিত সূর্যোদয়কালীন উদয় ও অস্তের মাধ্যমে অধিমাস যোগের পদ্ধতি ধীরে ধীরে ফসল দেখে করা অনিয়মিত অভিজ্ঞতাভিত্তিক পদ্ধতির জায়গা নিয়েছিল নিশ্চয়ই। এর একটি সুস্পষ্ট ইঙ্গিত পাওয়া যায় একটি তালিকায় (অজানা কোনো শতাব্দীর একটি মূল লেখার আসিরীয় অনুলিপি)। এতে নক্ষত্র বা তারামণ্ডলের ৩৬টি নাম আছে, প্রতি মাসের জন্য তিনটি করে, যার প্রতিটির প্রথমটি স্পষ্টতই সেই মাসের সূর্যোদয়কালীন উদয়ের সঙ্গে সম্পর্কিত। নিসান্নু মাসের জন্য পাই দিলগান (অর্থাৎ মেষ ও তিমি)। আইয়ারুর জন্য প্রথম নামটি মুলমুল (কৃত্তিকা; "মুল" মানে "নক্ষত্র")। সিমান্নুর জন্য সিবজিয়ান্নু (কালপুরুষ), এবং এভাবে বাকিগুলো। এ ছাড়া আরেকটি অনেকটা ক্ষতিগ্রস্ত লেখায় লিপিবদ্ধ আছে: "দিলগান নক্ষত্র নিসান্নু মাসে দেখা দেয়; নক্ষত্রটি অনুপস্থিত থাকলে মাসটিকে অবশ্যই… মুলমুল নক্ষত্র আইয়ারু মাসে দেখা দেয়…" নক্ষত্রের আবির্ভাব বা আবির্ভাবে দেরি স্পষ্টতই নির্দিষ্ট কিছু কাজের নির্দেশনা দেয়।


‌% ch4.tex
\chapter{আসিরিয়ান জ্যোতিষশাস্ত্র}

\epigraph{``\textit{বল মহাবিশ্বের মহাকাশ ফাড়ি’\\
চন্দ্র সূর্য গ্রহ তারা ছাড়ি’\\
ভূলোক দ্যুলোক গোলক ভেদিয়া\\
খোদার আসন ‘আরশ’ ছেদিয়া,\\
উঠিয়াছি চির-বিস্ময় আমি বিশ্ববিধাতৃর}''}{\textit{নজরুল}}


খ্রিস্টপূর্ব প্রথম সহস্রাব্দের শুরুতে উত্তরে, টাইগ্রিসের উজানে, একটি নতুন সেমিটিক শক্তির উত্থান ঘটে। এর রাজধানীর নাম ছিল আশুর, যা প্রধান দেবতারও নাম। ব্যাবিলনসহ আশপাশের রাষ্ট্রগুলোর সঙ্গে অবিরাম সংঘাতের মধ্য দিয়ে আশুর বিস্তার লাভ করে এবং মেসোপটেমীয় সমভূমির শাসকশক্তিতে পরিণত হয়। খ্রিস্টপূর্ব ৮০০ সালের দিকে আসিরিয়া নিকট-এশিয়ার সবচেয়ে শক্তিশালী রাষ্ট্র হয়ে ওঠে। এর রাজারা ছিলেন তিগলাথ-পিলেসার (খ্রিস্টপূর্ব ৭৪৫–৭২৭), শালমানেসার (৭২৬–৭২২), সারগন (৭২২–৭০৫), সেনাখেরিব (৭০৫–৬৮২), এসারহাদ্দোন (৬৮২–৬৬৮) এবং আশুরবানিপাল (খ্রিস্টপূর্ব ৬৬৮–৬২৬)। তাঁদের আমলে একের পর এক বড় এবং প্রায়ই নিষ্ঠুর যুদ্ধে সিরিয়া, প্যালেস্টাইন, ফিনিশিয়া, এমনকি কখনো কখনো মিশরও জয় করা হয়। সীমানা এশিয়া মাইনর, আর্মেনিয়া ও মিডিয়ার দিকে বিস্তৃত হয়। মার্কিন অ্যাসিরিওলজিস্ট ওলমস্টেডের মতে, আসিরীয় বিজয়ের প্রধান বস্তুগত কারণ ছিল অস্ত্রে লোহার ব্যবহার।

নতুন রাজধানী নিনেভে ছিল একটি বিশাল সামরিক সাম্রাজ্যের রাজনৈতিক কেন্দ্র, আর সে হিসেবে একে অপূর্ব সব ভবন দিয়ে সাজানো হয়েছিল। কিন্তু ব্যাবিলন ছিল একটি বড় ও সমৃদ্ধ বাণিজ্যকেন্দ্র, যার ধনী নাগরিকেরা অনেকাংশে নিজেরাই নিজেদের শাসন করত। প্রাচীন সংস্কৃতির পূজনীয় পীঠস্থান হিসেবে ব্যাবিলন তার মর্যাদা ধরে রেখেছিল। আসিরীয় রাজারা এর গুরুত্ব স্বীকার করতেন। তাঁরা নিজেরা ব্যাবিলনে গিয়ে "মারদুকের হাত ধরতেন", অর্থাৎ আনুষ্ঠানিকভাবে শাসনভার গ্রহণ করতেন, অথবা কোনো আত্মীয়কে অধীনস্থ রাজা হিসেবে মনোনীত করতেন।

রুক্ষ ও যুদ্ধপ্রিয় বিজেতারা প্রায়ই যা করে, আসিরীয়রাও তা-ই করল: তারা বিজিতদের সংস্কৃতি গ্রহণ করল এবং তার রূপ ও ঐতিহ্য বহন করে চলল। আসিরীয় দেবমণ্ডলী ব্যাবিলনেরটির মতোই ছিল, পার্থক্য কেবল এই যে এখন আশুর হলেন দেবতাদের নতুন দেবতা। একই পুরোহিতসমাজ নতুন শাসকদের সেবায় পুরোনো পবিত্র রীতিতে চিন্তা করত, কাজ করত ও লিখত, এবং আগের মতোই একই সামাজিক দায়িত্ব পালন করত। ভাস্কর্যকলার নতুন করে বিকাশ ঘটল। এর এক কারণ, কারিগরেরা ছিল নতুন প্রাণশক্তিসম্পন্ন একটি নতুন জাতি, আর নতুন শাসকেরা বিজয় ও লুণ্ঠনের মাধ্যমে ধনী হয়েছিলেন। আরেক কারণ, নিম্নভূমির ইটের বদলে আশপাশের পাহাড়ের পাথর নির্মাণসামগ্রী হিসেবে ব্যবহৃত হতো। সুন্দর উদ্‌গত খোদাইয়ের জন্য এটি ছিল চমৎকার উপকরণ, আর সেই খোদাইগুলো আজ পশ্চিমের জাদুঘরগুলোর শোভা বাড়াচ্ছে। এই বিকাশ চূড়ায় পৌঁছায় যখন সম্পদ ও সংস্কৃতি বাড়ার সঙ্গে সঙ্গে প্রথম দিকের কঠোর যোদ্ধাদের জায়গায় এলেন এমন রাজপুত্রেরা, যাঁরা শিল্প ও বিজ্ঞান ভালোবাসতেন এবং তার পৃষ্ঠপোষকতা করতেন। তখনই আশুরবানিপাল তাঁর প্রাসাদে একটি গ্রন্থাগার স্থাপন করেন। তিনি ব্যাবিলোনিয়ার সব প্রাচীন স্থান ও মন্দির থেকে পুরোনো লেখাগুলো সংগ্রহ করে অনুলিপি করার আদেশ দেন। হাজার হাজার মাটির ফলক সুশৃঙ্খল সারিতে রাখা হয়েছিল, পাশে তাদের শিরোনাম লেখা থাকত। ক্যাটালগ, অভিধান ও ভাষ্য দিয়ে এগুলোর পরিপূরক ও ব্যাখ্যা দেওয়া হতো। নতুন নতুন মহাফেজখানার দলিল, রাজার কাছে পাঠানো প্রতিবেদন এবং কর্মকর্তাদের সঙ্গে তাঁর চিঠিপত্র দিয়ে সংগ্রহটি অবিরাম বাড়ানো হতো। কেবল এই স্থান থেকেই খুঁড়ে পাওয়া ১৩,০০০-এর বেশি খণ্ড ব্রিটিশ মিউজিয়ামে সংরক্ষিত আছে, আর অন্যান্য স্থান থেকে আরও বহু হাজার পাওয়া গেছে। এগুলো থেকে সেই সমাজের রীতিনীতি ও ধ্যানধারণা, ব্যবসায়িক জীবন ও সংস্কৃতি, ধর্ম ও অর্থনীতি, এবং তার জ্যোতির্বিজ্ঞানেরও একটি ভালো চিত্র পাওয়া যায়।

অসংখ্য লেখায় আগের যুগের পঞ্জিকাই পাওয়া যায়। চান্দ্র মাস শুরু হতো বাঁকা চাঁদের প্রথম আবির্ভাবের সন্ধ্যা থেকে, আর ১২টি মাসের সঙ্গে মাঝে মাঝে একটি ত্রয়োদশ মাস যোগ করে বছর পূর্ণ করা হতো। স্পষ্ট উল্লেখ না থাকলেও এটি প্রায় নিশ্চিত যে নক্ষত্রসংক্রান্ত ঘটনা দিয়ে, সাধারণত নক্ষত্রের প্রভাতী উদয় দিয়ে, অধিমাস যোগ নিয়ন্ত্রণ করা হতো। আগে উল্লিখিত ৩ × ১২টি মাস-নক্ষত্রের তালিকার একটি অনুলিপি আশুরবানিপালের গ্রন্থাগারে আছে, যা এর একটি ইঙ্গিত হতে পারে। জর্জ স্মিথের প্রকাশিত একটি লেখায় আরেকটি পদ্ধতির ইঙ্গিত পাওয়া যায়: "নিসান্নুর প্রথম দিনে চাঁদ ও মুলমুল নক্ষত্র [কৃত্তিকা] একসঙ্গে থাকলে বছরটি সাধারণ; নিসান্নুর তৃতীয় দিনে চাঁদ ও মুলমুল নক্ষত্র একসঙ্গে থাকলে বছরটি পূর্ণ।" বাক্যের শেষ অংশের অর্থ হলো, কৃত্তিকা সূর্যাস্তের অনেক পরেও দৃশ্যমান থাকে। বসন্তের এত আগে এমন হলে একটি ত্রয়োদশ মাস যোগ করা প্রয়োজন।

নক্ষত্র ও তারামণ্ডলগুলোকে তিনটি ভাগে ভাগ করা হয়েছিল: উত্তরের, মধ্যের ও দক্ষিণের তারামণ্ডল। এগুলো যথাক্রমে এনলিল, আনু ও এয়ার এলাকা নামে পরিচিত ছিল।

তবে নক্ষত্র পর্যবেক্ষণের প্রধান উদ্দেশ্য আর পঞ্জিকা ছিল না। আসিরীয় যুগে মানুষের চিন্তা ও চর্চাকে সবচেয়ে জোরালোভাবে নির্ধারণ করত জ্যোতিষশাস্ত্র, অর্থাৎ এই ধারণা যে নক্ষত্রের গতিপথ পৃথিবীর ঘটনাবলির জন্য তাৎপর্যপূর্ণ। এখন আকাশের ঘটনাবলি অনেক গভীর আগ্রহ নিয়ে অধ্যয়ন করা হতো মানুষের ভাগ্য সম্পর্কে, বিশেষ করে রাজা ও সাম্রাজ্যের ভাগ্য সম্পর্কে, শুভাশুভ লক্ষণ খুঁজতে।

শুভাশুভ লক্ষণে এই বিশ্বাস আদিম মানুষের মধ্যেও ছিল। এটি ছিল তার এই বিশ্বাসের স্বাভাবিক পরিণতি যে সে অদৃশ্য আত্মিক শক্তি দিয়ে পরিবেষ্টিত, আর সেই শক্তিগুলো তার জীবন ও কাজকে প্রভাবিত করে। তাদের অনুগ্রহ ও সাহায্য লাভ করা, তাদের শত্রুতা প্রশমিত বা প্রতিহত করা এবং তাদের অভিপ্রায় জানা ছিল তার কাছে জীবন-মরণের প্রশ্ন। ভূত তাড়ানো, নৈবেদ্য, মন্ত্র, তাবিজ এবং কাজের সঙ্গে জড়িত জাদু তার দৈনন্দিন জীবন জুড়ে থাকত। এই আত্মাদের বেশিরভাগের বাস ছিল আকাশে। তাই প্রাচীন কালেই মেসোপটেমীয় পুরোহিতদের মনে নক্ষত্র ও মানুষের ভাগ্যের মধ্যে ঘনিষ্ঠ সম্পর্কের ধারণা জন্মেছিল। স্বর্গ খুব দূরে ছিল না; পৃথিবীর শাসক ও তাদের প্রজারা ওপরের দেবতাদের কাছাকাছিই ছিল। লাগাশের রাজা গুদেয়া (আনুমানিক খ্রিস্টপূর্ব ২৫০০) একটি পাথরের সিলিন্ডারে খোদাই করা লেখায় একটি মন্দির নির্মাণের বর্ণনা দিয়েছেন। তিনি বলেছেন, কীভাবে স্বপ্নে এয়ার কন্যা দেবী নিসাবা তাঁর সামনে আবির্ভূত হয়েছিলেন: "তাঁর হাতে ছিল উজ্জ্বল লেখনী; তিনি শুভ আকাশ-লক্ষণযুক্ত একটি ফলক বহন করছিলেন এবং চিন্তামগ্ন ছিলেন।" আরও পরে: "তিনি মন্দির নির্মাণের জন্য শুভ নক্ষত্রের কথা ঘোষণা করলেন।" প্রথম ব্যাবিলনীয় সাম্রাজ্যের নিন-দার-আন্না লেখায় গ্রহসংক্রান্ত প্রতিটি ঘটনার পর একটি করে লক্ষণ পাওয়া যায়। এতে কোনো সন্দেহ নেই যে এই ভবিষ্যদ্বাণীগুলো জোগানোর জন্যই জ্যোতির্বৈজ্ঞানিক তথ্য সংগ্রহ করা হতো।

ব্যাবিলন ও অন্যান্য নগরের মন্দিরের পুরোহিতদের মধ্যে এমন বিশ্বাস ছিল একটি জীবন্ত শক্তি। ব্যাবিলনীয় সংস্কৃতির প্রভাবাধীন সব অঞ্চলে, আসিরিয়াসহ, এটি ছড়িয়ে পড়ে। বিজয়ের মাধ্যমে সাম্রাজ্য বিস্তারে উদ্‌গ্রীব শক্তিশালী রাজাদের দরবারে ভবিষ্যৎ আগে থেকে জানার প্রয়োজন ছিল বিশেষভাবে বেশি, আর সেখানে জ্যোতিষশাস্ত্র উর্বর ক্ষেত্র খুঁজে পেল। রাজদরবারের জ্যোতিষীদের প্রতিটি বড় উদ্যোগের জন্য লক্ষণ খুঁজে বের করতে হতো। সব গুরুত্বপূর্ণ মন্দির থেকে রাজা নিয়মিত প্রতিবেদন পেতেন যে আকাশে কী ঘটছে এবং তার ব্যাখ্যা কী। এই প্রতিবেদনগুলো সেই মহাফেজখানায় সংরক্ষিত ছিল যাকে আমরা আশুরবানিপালের গ্রন্থাগার বলি, আর ব্যাখ্যার সুবিধার জন্য সব পুরোনো তথ্যের অনুলিপি সেখানে জড়ো করা হয়েছিল।

অবশ্য সব লক্ষণ নক্ষত্র থেকে আসত না। লক্ষণ সর্বত্রই পাওয়া যেত: বলির পশুর যকৃতে, যা পরিবর্তনশীল বলে লক্ষণের প্রচুর উৎস ছিল; পাখির ওড়ায়; গর্ভপাত, ভূমিকম্প, মেঘ, রংধনু ও জ্যোতির্বলয়ে। এমন সব নির্দেশিকা ছিল, যাতে প্রতিটি ঘটনা লিপিবদ্ধ থাকত, আর সেই সঙ্গে থাকত কিংবদন্তির রাজা আগাদের সারগনের কোন কাজ বা অভিজ্ঞতার সঙ্গে এর কোনোভাবে সম্পর্ক আছে।

স্বাভাবিকভাবেই এগুলোর মধ্যে জ্যোতির্বৈজ্ঞানিক ঘটনা বড় জায়গা জুড়ে ছিল। যেসব জ্যোতিষ্কের ঘটনায় বৈচিত্র্য ও অনিয়ম বেশি, সেগুলোই পূর্বাভাসের জন্য সবচেয়ে উপযোগী ছিল। তাই সবচেয়ে গুরুত্বপূর্ণ ছিল সূর্য ও স্থির নক্ষত্র নয়, বরং চাঁদ ও গ্রহগুলো। স্থির নক্ষত্রগুলো ঘটনাবলির জন্য অপরিবর্তনীয় পটভূমি হিসেবে কাজ করত। ন্যায়বিচারের সর্বদ্রষ্টা রক্ষক সূর্যদেব শামাশ দিনের পর দিন আকাশজুড়ে তাঁর অপরিবর্তনীয় পথে চলেন। বড়জোর তিনি স্বাভাবিকের চেয়ে নিষ্প্রভ বা বেশি লাল হতে পারতেন, কিংবা মেঘ বা গ্রহণের কারণে অদৃশ্য হতে পারতেন। অন্যদিকে চন্দ্রদেব সিন বারবার ফিরে আসা নানা রকম ঘটনার প্রাচুর্য দেখান। গ্রহগুলোর ক্ষেত্রেও একই কথা খাটে। তারা নক্ষত্রদের মাঝে সবচেয়ে অপ্রত্যাশিত উপায়ে ঘুরে বেড়ায়। কখনো স্থির দাঁড়িয়ে থাকে বা উল্টো পথে চলে, কখনো নিজেদের মধ্যে বা উজ্জ্বল নক্ষত্রগুলোর সঙ্গে মিলে সদা পরিবর্তনশীল বিন্যাস তৈরি করে। এদের মনে হতো তারাভরা আকাশের প্রান্তরে স্বেচ্ছায় ঘুরে বেড়ানো জীবন্ত সত্তা, আর ক্রমশ এরাই ব্যাবিলনীয় পুরোহিতদের মনোযোগের প্রধান বিষয় হয়ে উঠল। এরা ছিল সেই মহান দেবতাদের নক্ষত্র, যাঁরা জগৎ শাসন করতেন এবং এই উজ্জ্বল জ্যোতিষ্কগুলোর মধ্য দিয়ে নিজেদের প্রকাশ করতেন।

গ্রহগুলোর নামের আগে বসানো নির্ধারক চিহ্নটি সাধারণত ছিল "নক্ষত্র"-এর চিহ্ন, এখন প্রায়ই তা হয়ে গেল "দেবতা"-র চিহ্ন। শুক্র (যার নাম দিল-বাত) ছিল ইশতারের নক্ষত্র। বৃহস্পতি (উমানপা-উদ্দা, পরে বেশিরভাগ সময় সাগমেগার) ছিল মারদুকের নক্ষত্র। শনি ছিল নিনিবের নক্ষত্র, যিনিও একজন সৌর দেবতা। আর মঙ্গল ছিল মহামারির দেবতা নের্গালের নক্ষত্র। লাল মঙ্গলকে অশুভ নক্ষত্র মনে করা হতো, বৃহস্পতিকে শুভ নক্ষত্র। কিন্তু কার জন্য শুভ বা অশুভ? তা নির্ভর করত সময় ও স্থানের ওপর।

চারটি দেশের প্রতিটির জন্য আলাদা আলাদা মাস, আলাদা দিক এবং আলাদা তারামণ্ডল নির্দিষ্ট ছিল: আক্কাদ (ব্যাবিলোনিয়া), এলাম (পূর্বের পাহাড়ি অঞ্চল), আমুররু (পশ্চিমের মরুভূমি, পরে সিরিয়া) এবং সুবারতু (উত্তর)। একজন আসিরীয় জ্যোতিষী বলেছিলেন, "আমরাই সুবারতু।" বিভিন্ন তারামণ্ডলে গ্রহগুলোর আবির্ভাব, তাদের অবস্থানের স্থায়িত্ব বা দ্রুত গতি, একে অপরের সঙ্গে বা চাঁদের সঙ্গে তাদের মিলন ঘটনার একেবারে অন্তহীন বৈচিত্র্য তৈরি করত। এতে জ্যোতিষীরা ঐতিহ্য ও উদ্ভাবনী বুদ্ধি খাটানোর বিস্তৃত সুযোগ পেতেন। ব্যক্তিগত সিদ্ধান্ত টানার সময় মেলানোর ক্ষমতা ও কল্পনা কাজে লাগানোর, আবার বিচক্ষণভাবে চাটুকারিতা করারও সুযোগ থাকত।

দৃষ্টান্ত হিসেবে এখানে কয়েকটি লেখা দেওয়া হলো। এগুলো অ্যাসিরিওলজির শিক্ষার্থীদের জন্য আর. সি. টমসনের প্রকাশিত সংকলন থেকে নেওয়া।

মারদুকের নক্ষত্র বছরের শুরুতে দেখা দিলে সেই বছর শস্য সমৃদ্ধ হবে। বুধ নিসান্নুতে দেখা দিয়েছে। কোনো গ্রহ [বুধ] লি নক্ষত্রের [অ্যালডেবারান] কাছে এলে এলামের রাজা মারা যাবেন।… বুধ বৃষ রাশিতে দেখা দিয়েছে; এটি শুগি [কৃত্তিকা, পারসিউস] পর্যন্ত নেমে এসেছিল [?]।

শুক্র পশ্চিমে অদৃশ্য হয়। আবু মাসে শুক্র নিষ্প্রভ হয়ে অদৃশ্য হলে এলামে হত্যাযজ্ঞ হবে। আবু মাসে প্রথম থেকে ত্রিশতম দিন পর্যন্ত শুক্র দেখা দিলে বৃষ্টি হবে, আর দেশের ফসল সমৃদ্ধ হবে। মাসের মাঝামাঝি শুক্র পূর্বে সিংহ রাশিতে দেখা দিয়েছে। নিরগাল-ইতির থেকে।

শুক্র যখন তার অবস্থান স্থির করে, তখন রাজার আয়ু দীর্ঘ হবে, দেশে ন্যায়বিচার থাকবে। শুক্র যখন এয়ার পথে থাকে…

দুজু মাসে মঙ্গল দৃশ্যমান; এটি নিষ্প্রভ।… মঙ্গল অস্পষ্টভাবে [?] মধ্যগগনে পৌঁছে উজ্জ্বল হয়ে উঠলে এলামের রাজা মারা যাবেন। দেবতা নের্গাল অদৃশ্য হওয়ার সময় যখন ছোট হয়ে আসেন এবং আকাশের নক্ষত্রদের মতো খুব অস্পষ্ট হন, তখন তিনি আক্কাদের প্রতি দয়া দেখাবেন।… মঙ্গল নিষ্প্রভ হলে শুভ; উজ্জ্বল হলে অশুভ। মঙ্গল বৃহস্পতিকে অনুসরণ করলে সেই বছর শুভ হবে। বুল্লুতু থেকে। 
মঙ্গল আল্লুলের [কর্কট] এলাকায় প্রবেশ করেছে। এটি লক্ষণ হিসেবে গণ্য হয় না। এটি থামেনি, অপেক্ষা করেনি, বিশ্রাম নেয়নি; দ্রুত এগিয়ে গেছে। বিল-না-সির থেকে।

বৃহস্পতি পশ্চিমে চলে গেলে নিরাপদ বসবাস হবে, দেশে কোমল শান্তি নেমে আসবে। এটি আল্লুলের সামনে দেখা দিয়েছে। বৃহস্পতি বেলের পথে উজ্জ্বলতা ধারণ করে নিবিরু (হয়ে উঠলে?), আক্কাদ প্রাচুর্যে উপচে পড়বে, আক্কাদের রাজা শক্তিশালী হবেন।… আগুনের মতো একটি বড় নক্ষত্র পূর্ব থেকে উঠে পশ্চিমে অদৃশ্য হলে যুদ্ধে শত্রুসৈন্য নিহত হবে। রাজত্বের শুরুতে বৃহস্পতিকে সঠিক অবস্থানে দেখা গেছে; দেবতাদের প্রভু তোমাকে সুখী করুন এবং তোমার আয়ু দীর্ঘ করুন। দামকার পুত্র আশারিদু থেকে।

এখানে উল্লিখিত অগ্নিময় নক্ষত্রটি সম্ভবত একটি উল্কা ছিল। পরের লেখাটিতেও একই ব্যাপার:

রাতের এক কাস-বু [দুই ঘণ্টা] পার হওয়ার পর একটি বড় নক্ষত্র উত্তর থেকে দক্ষিণে জ্বলে উঠল। এর লক্ষণ রাজার আকাঙ্ক্ষার অনুকূল। আক্কাদের রাজা তাঁর উদ্দেশ্য সফল করবেন। রাজার সেবক জ্যেষ্ঠ আশারিদু থেকে।

রাজা এসারহাদ্দোনকে মার-ইশতারের লেখা একটি চিঠিতে (খ্রিস্টপূর্ব ৬৬৮) আমরা পড়ি:

…প্রথম মাসে… ২৯ তারিখে বৃহস্পতিকে সরিয়ে নেওয়া হলো… এখন সে এক মাস পাঁচ দিন আকাশে অনুপস্থিত থেকেছে [অর্থাৎ অদৃশ্য ছিল]; তৃতীয় মাসের ষষ্ঠ দিনে বৃহস্পতি কালপুরুষ অঞ্চলে দৃশ্যমান হলো; তার নির্ধারিত সময়ের চেয়ে পাঁচ দিন বেশি সে পার করছিল। এটি এই লক্ষণের সঙ্গে মেলে: বৃহস্পতি তৃতীয় মাসে দেখা দিলে দেশ বিধ্বস্ত হবে এবং শস্য দুর্মূল্য হবে… বৃহস্পতি কালপুরুষে প্রবেশ করলে দেবতারা দেশকে গ্রাস করবেন।

এই এবং এ জাতীয় অন্যান্য লেখায় (যেখানে কিছু অভিব্যক্তির অর্থ এখনো অনিশ্চিত) আমরা এমন জ্যোতির্বৈজ্ঞানিক পর্যবেক্ষণ পাই, যা আগে পঞ্জিকা বা দিক নির্ণয়ের জন্য যা দরকার হতো তার চেয়ে অনেক বেশি বৈচিত্র্যময় ও বিস্তারিত। জ্যোতিষশাস্ত্র মানুষের জীবনকে আকাশের সঙ্গে এত ঘনিষ্ঠভাবে যুক্ত করেছিল যে নক্ষত্র ও তাদের পরিভ্রমণ তার চিন্তা ও কর্মকাণ্ডে গুরুত্বপূর্ণ স্থান নিতে শুরু করল। উজ্জ্বল দেবতারা সেই আশ্চর্য কক্ষপথগুলোতে তার ভাগ্যই বুনে চলতেন। এখন যেহেতু তার চোখ খুলে গিয়েছিল এবং আগ্রহ জেগেছিল, মানুষ গ্রহগুলোর গতিপথের সঙ্গে ক্রমশ আরও পরিচিত হয়ে উঠল। প্রতিবেদনগুলোতে আমরা তা দেখেছি: বৃহস্পতি পশ্চিমে যায়; মঙ্গল বৃশ্চিকে দাঁড়িয়ে ছিল, ঘুরে কম উজ্জ্বলতা নিয়ে এগিয়ে যায়; শুক্র তার জায়গায় থাকে; আর শনি সিংহে দেখা দিয়েছে। পর্যায়কালের কোনো তত্ত্ব ছিল না (এই সব বিবৃতিতে বছর ও তারিখ অনুপস্থিত)। তবু পর্যবেক্ষকদের মনে নিশ্চয়ই একটি পরিচিত ও প্রত্যাশিত নিয়মিততার ধারণা জন্মেছিল। গ্রহটি সঠিক জায়গায় ছিল বা তার সময় পেরিয়ে গেছে, এমন বিবৃতিতে তা প্রকাশ পায়।

গ্রহের পরেই জ্যোতিষীর কাছে প্রধান জ্যোতিষ্ক ছিল চাঁদ। এর কালগণনাগত তাৎপর্য এখন জ্যোতিষীয় আগ্রহের কাছে জায়গা ছেড়ে দিতে বাধ্য হলো। কালগণনার ত্রুটিগুলোও জ্যোতিষীদের কাছে ভবিষ্যদ্বাণীর উপাদান হয়ে উঠল। চাঁদ সময়মতো না এলে তা বেশিরভাগ ক্ষেত্রে ছিল অশুভ লক্ষণ। মেঘের কারণে চাঁদ অদৃশ্য থাকায়, কিংবা হয়তো দুর্যোগের সময়ে অবহেলার কারণে, নতুন বাঁকা চাঁদ মাসের ২৮ বা ২৯ তারিখেও দেখা দিতে পারত। একইভাবে পূর্ণিমা, যা স্বাভাবিকভাবে ১৪ তারিখে হওয়ার কথা, ১৩, ১৫ বা ১৬ তারিখে হতে পারত।

চাঁদ প্রথম দিনে দেখা দিলে নীরবতা থাকবে, দেশ পরিতৃপ্ত থাকবে।…

নিসান্নুর ত্রিশ তারিখে চাঁদ দেখা দিলে সুবারতু আখলামুকে গ্রাস করবে; আমুররুতে একটি বিদেশি ভাষা প্রাধান্য পাবে। আমরাই সুবারতু। ত্রিশতম দিনে চাঁদ দেখা দিলে দেশে শীত পড়বে। তেবিতুর ১৪ তারিখে চাঁদ সূর্য ছাড়া দেখা দিয়েছে; শাবাতুতে চাঁদ দিন পূর্ণ করবে।…

ত্রয়োদশ দিনে চাঁদ ও সূর্যকে একসঙ্গে দেখা গেলে নীরবতা থাকবে না; দেশে ব্যবসা-বাণিজ্য ব্যর্থ হবে, শত্রু দেশ দখল করবে। আপলা থেকে।

চাঁদ সূর্যের কাছে পৌঁছে তার সঙ্গে দৃষ্টির আড়ালে মিলিয়ে গেলে… দেশে সত্য প্রতিষ্ঠিত হবে এবং পুত্র পিতার সঙ্গে সত্য কথা বলবে। ১৪ তারিখে দেবতাকে দেবতার সঙ্গে দেখা গেছে।… ১৪ তারিখে চাঁদ ও সূর্যকে একসঙ্গে দেখা গেলে নীরবতা থাকবে, দেশ পরিতৃপ্ত থাকবে; দেবতারা আক্কাদের সুখ চান।…

চাঁদ সূর্যের জন্য অপেক্ষা না করে অদৃশ্য হলে সিংহ ও নেকড়ের তাণ্ডব হবে।… ১৫ তারিখে একে সূর্যের সঙ্গে দেখা গেছে; পরে তিশরিতুতে চাঁদ দিন পূর্ণ করবে।… বালাসি থেকে।

প্রথম দিনে আমি রাজার কাছে এভাবে পাঠিয়েছিলাম: ১৪ তারিখে চাঁদকে সূর্যের সঙ্গে দেখা যাবে।… ১৪ তারিখে চাঁদকে সূর্যের সঙ্গে দেখা গেছে।

ষোড়শ দিনে চাঁদ ও সূর্যকে একসঙ্গে দেখা গেলে রাজা রাজার প্রতি শত্রুতা পাঠাবেন। রাজা এক মাস ধরে তাঁর প্রাসাদে অবরুদ্ধ থাকবেন। শত্রুর পা দেশের বিরুদ্ধে থাকবে; শত্রু বিজয়গর্বে তার দেশে কুচকাওয়াজ করবে। দুজুর ১৪ বা ১৫ তারিখে চাঁদকে সূর্যের সঙ্গে দেখা না গেলে রাজা তাঁর প্রাসাদে অবরুদ্ধ হবেন। ষোড়শ দিনে দেখা গেলে তা সুবারতুর জন্য শুভ, আক্কাদ ও আমুররুর জন্য অশুভ। আকেল্লানু থেকে।

এমন পর্যবেক্ষণ ও সিদ্ধান্ত বুঝতে হলে পূর্ণিমার কাছাকাছি সময়ে কোন ঘটনাগুলো দেখা যায় তা বিবেচনা করতে হবে। মাস যদি স্বাভাবিকভাবে শুরু হয়, অর্থাৎ যে সন্ধ্যায় বাঁকা চাঁদ প্রথমবার দৃশ্যমান হয় (সাধারণত সংযোগের ১.৪ দিন পর) সেদিন থেকেই প্রথম দিন শুরু হয়, তাহলে পূর্ণিমা পড়বে ১৪ তারিখের রাতে। কারণ পূর্ণিমা আসে সংযোগের ১৪.৭ দিন পর (গড়ে; চাঁদের গতির অনিয়মের কারণে এটি ½ দিন কম বা বেশি হতে পারে)। ১৩ তারিখের রাতে চাঁদ তখনো পূর্ণ হয়নি। সূর্যাস্তের সময়েই চাঁদ পূর্বে দৃশ্যমান থাকে, আর সূর্যোদয়ের আগেই সে অস্ত যায়; সে সূর্যের জন্য অপেক্ষা করে না। ১৪ তারিখের রাতেও সূর্যাস্তের সময় চাঁদ দৃশ্যমান থাকে, কিন্তু সূর্য ওঠার সময়ও সে অস্ত যায়নি, পশ্চিমে নিষ্প্রভ ও ম্লান হয়ে আছে। আসিরীয় জ্যোতিষীরা একে প্রকাশ করতেন এভাবে: "চাঁদকে সূর্যের সঙ্গে দেখা গেছে", বা "দেবতা [সিন] দেবতা [শামাশ]-এর সঙ্গে দেখা দিয়েছেন", বা "চাঁদ সূর্যের কাছে পৌঁছায়"। এমন স্বাভাবিক অবস্থার সঙ্গে নানা রকম শুভ লক্ষণ জড়িত ছিল। কিন্তু যদি "চাঁদ সূর্যের জন্য অপেক্ষা না করে" সূর্যোদয়ের আগেই অস্ত যায়, তাহলে পূর্ণিমা আসবে পরে। তখন পরের রাত, অর্থাৎ ১৫ তারিখের আগে প্রভাতে চাঁদকে সূর্যের সঙ্গে দেখা যাবে না।

এই পর্যবেক্ষণগুলোতে জ্যোতিষশাস্ত্র কালগণনার সঙ্গে মিশে গিয়েছিল। পুরোহিত-জ্যোতির্বিদেরা লক্ষ করতেন যে আকাশে সব ঠিকঠাক আছে, তাই পৃথিবীতে শান্তি থাকবে। অথবা তাঁরা পঞ্জিকায় কোনো বিশৃঙ্খলা লক্ষ করতেন, যা অশুভ লক্ষণ ছাড়াও নির্দেশ করত যে কিছু একটা সংশোধন করতে হবে। প্রথম আবির্ভাবের দিন থেকে হিসাব করলে পূর্ণিমা ১৪ তারিখে প্রত্যাশিত ছিল, কিন্তু অনিয়মের কারণে তা না-ও হতে পারত। পূর্ণিমা ১৩ তারিখে এলে তাঁরা সিদ্ধান্ত নিতেন যে মাসটি হয়তো মাত্র ২৯ দিনের হবে। ১৪ তারিখের পরে চাঁদকে সূর্যের সঙ্গে দেখা গেলে মাসটিকে ৩০ দিনের হতে হতো; তখন সে "দিন পূর্ণ করত"। ১৪ তারিখে একসঙ্গে দেখা গেলে এবং আগের মাস ৩০ দিনের হলে পরের মাসটি হতো মাত্র ২৯ দিনের।

আদারুতে চাঁদ দিন পূর্ণ করেছে; ১৪ তারিখে চাঁদকে সূর্যের সঙ্গে দেখা যাবে; নিসান্নুতে চাঁদ দিন ফিরিয়ে নেবে।…

এভাবে জ্যোতিষীয় উদ্দেশ্য চাঁদ পর্যবেক্ষণে বেশি মনোযোগ দাবি করে কালগণনার জন্য আরও ভালো ভিত্তি জুগিয়েছিল। শৃঙ্খলা ও নিয়মিততা উপলব্ধি করা হলো এবং ভবিষ্যদ্বাণীর কাজে লাগানো হলো।

অবশ্যই চাঁদে আরও অনেক ঘটনা পর্যবেক্ষণ করা যেত: এর রং ও উজ্জ্বলতা, শিংয়ের আকৃতি, পৃথিবী-আলো (চাঁদ একটি "আগু", অর্থাৎ মুকুট বা রাজকীয় টুপি পরে আছে), একটি কুয়াশাময় ছটা, বা একটি বলয় (জ্যোতির্বলয়)। এমন বলয়কে প্রায়ই মনে করা হতো ভেড়ার চারণভূমি ঘিরে রাখা একটি বেড়া, যার কেন্দ্রে রাখাল হিসেবে আছে চাঁদ। গ্রহের ব্যাবিলনীয় চিহ্ন "লুবাত"-এর অর্থ "পথহারা ভেড়া"। অথবা বলয়টিকে একটি নদী হিসেবে দেখা হতো; কিংবা এর অর্থ ছিল অবরোধ, আর বলয়ের ভেতরের গ্রহ বা নক্ষত্রগুলো নির্দেশ করত কে অবরুদ্ধ। বলয়টি বন্ধ না হলে স্বাভাবিকভাবেই তার অর্থ ছিল পলায়ন।

জ্যোতির্বলয় চাঁদকে ঘিরে থাকলে এবং বৃহস্পতি তার ভেতরে থাকলে রাজা অবরুদ্ধ হবেন। বলয়টি বিচ্ছিন্ন ছিল; এটি অশুভ নির্দেশ করে না।… নাবু-শুমা-ইশকুন থেকে।

জ্যোতির্বলয় চাঁদকে ঘিরে থাকলে এবং সুদুন তার ভেতরে থাকলে একজন রাজা মারা যাবেন এবং তাঁর দেশ সংকুচিত হবে; এলামের রাজা মারা যাবেন। সুদুন হলো মঙ্গল, মঙ্গল আমুররুর নক্ষত্র; এটি আমুররু ও এলামের জন্য অশুভ। শনি আক্কাদের নক্ষত্র, এটি রাজা, আমার প্রভুর জন্য শুভ। রাজার সেবক ইরাশ্‌শি-ইলু থেকে।

গ্রহণকেও অত্যন্ত গুরুত্বপূর্ণ লক্ষণ মনে করা হতো। সব যুগেই এর অপ্রত্যাশিত আবির্ভাব মানুষকে গভীরভাবে নাড়া দিয়েছে। নিয়মিত ঘটনাগুলোকেই যেখানে এত তাৎপর্য দেওয়া হতো, সেখানে এর প্রভাব আরও বেশি ছিল। গ্রহণ থেকে বহু রকম ভবিষ্যদ্বাণী করা হতো। সূর্যগ্রহণের ক্ষেত্রে মাস ও আকাশে অবস্থান লিখে রাখা হতো। সূর্য যখন "চাঁদের আকার ধারণ করত", তখন শিংগুলোর চেহারা ও দিকও লিখে রাখা হতো। মাসের ২৭ বা ২৮ তারিখ ছাড়া অন্য তারিখে গ্রহণের বিবরণ পাওয়া গেলে বোঝা যায়, জ্যোতির্বৈজ্ঞানিক গ্রহণ ছাড়াও ধুলিঝড় সূর্যকে অন্ধকার করে থাকতে পারে।

চন্দ্রগ্রহণ থেকে পাওয়া লক্ষণ ছিল অনেক বেশি। চাঁদের পূর্ণগ্রহণ হলে তা সব দেশের জন্য লক্ষণ ছিল। আংশিক আচ্ছন্ন হলে এর চার দিকের প্রতিটি আলাদা একটি দেশের সঙ্গে সম্পর্কিত ছিল, আর মাস, দিন ও প্রহর (রাতের পাহারা) অনুযায়ী তাৎপর্য আলাদা হতো। তাই "এনুমা আনু এনলিল" নামের প্রাচীন লক্ষণের বিশাল সংকলনে এই নির্দেশনা দেওয়া হয়েছে:

চাঁদের গ্রহণ হলে তুমি নিখুঁতভাবে পর্যবেক্ষণ করবে মাস, দিন, রাতের প্রহর, বাতাস, গতিপথ এবং সেই নক্ষত্রগুলোর অবস্থান যাদের এলাকায় গ্রহণ ঘটছে। এর মাস, দিন, রাতের প্রহর, বাতাস, গতিপথ ও নক্ষত্র সম্পর্কিত লক্ষণগুলো তুমি নির্দেশ করবে।

তাই সিমান্নু মাসে রাতের শেষে ঘটে যাওয়া একটি গ্রহণের প্রতিবেদনে আমরা পাই:

প্রভাতী প্রহরে গ্রহণের অর্থ রোগ।… প্রভাতী প্রহর হলো এলাম, ১৪ তারিখ হলো এলাম, সিমান্নু হলো আমুররু, দ্বিতীয় দিক হলো আক্কাদ।… প্রভাতী প্রহরে গ্রহণ হয়ে প্রহরটি পূর্ণ করলে, এবং উত্তুরে বাতাস বইলে, আক্কাদের অসুস্থরা সেরে উঠবে। গ্রহণ প্রথম দিকে শুরু হয়ে দ্বিতীয় দিকে থামলে এলামে হত্যাযজ্ঞ হবে; গুতি আক্কাদের কাছে আসবে না।… গ্রহণ হয়ে দ্বিতীয় দিকে থামলে দেবতারা দেশের প্রতি দয়া দেখাবেন। সিমান্নুতে চাঁদ অন্ধকার হলে এক বছর পর রামানু [ঝড়ের দেবতা] প্লাবন ঘটাবেন। সিমান্নুতে চাঁদের গ্রহণ হলে বন্যা হবে, আর দেশের জলজ উৎপাদন প্রচুর হবে।…

লক্ষণের এক সমৃদ্ধ তালিকা; প্রতিটি খুঁটিনাটির তাৎপর্য ছিল। এত বিস্তারিত বিবরণ থাকলে আমরা গ্রহণের বেশ নির্ভুল পর্যবেক্ষণের কথা বলতে পারি। তাই আশ্চর্যের কিছু নেই যে বহু শতাব্দী পরে টলেমি চাঁদের পর্যায়কাল নির্ণয়ের জন্য যে প্রথম পর্যবেক্ষিত গ্রহণগুলো ব্যবহার করেছিলেন, সেগুলো এই যুগেরই: খ্রিস্টপূর্ব ৭২১ সালের ১৯ মার্চ, খ্রিস্টপূর্ব ৭২০ সালের ৮ মার্চ এবং খ্রিস্টপূর্ব ৭২০ সালের ১ সেপ্টেম্বর। এখানে কেবল প্রয়োজন ছিল এই তারিখগুলোর পর কত বছর ও কত দিন পেরিয়েছে তার নিখুঁত জ্ঞান। হতে পারে, এই তথ্যের অভাবেই আরও আগের গ্রহণগুলো ব্যবহার করা যায়নি। আর এ কারণেই টলেমি তাঁর রাজাদের তালিকা ও তাঁদের শাসনের ধারাবাহিক বছরগুলো শুরু করেছিলেন ব্যাবিলনীয় রাজা নাবোনাসারকে (খ্রিস্টপূর্ব ৭৪৭) দিয়ে।

ব্যাবিলনীয় প্রহরীরা নিজেরা নিশ্চয়ই এত যত্নের সঙ্গে পর্যবেক্ষিত ঘটনাবলিতে নিয়মিততা লক্ষ করেছিলেন। তাঁদের প্রতিবেদন দেখায় যে তাঁরা গ্রহণ প্রত্যাশা করতেন এবং ঘটনাগুলো আগাম ঘোষণা করতেন:

১৪ তারিখে একটি গ্রহণ ঘটবে; এটি এলাম ও আমুররুর জন্য অশুভ, রাজা, আমার প্রভুর জন্য শুভ; রাজা, আমার প্রভু, নিশ্চিন্তে থাকুন। এটি শুক্র ছাড়াই দেখা যাবে। রাজা, আমার প্রভুকে আমি বলছি: একটি গ্রহণ হবে। রাজার সেবক ইরাশ্‌শি-ইলু থেকে।

দেশসমূহের রাজা, আমার প্রভুর প্রতি, তোমার সেবক বিল-উসুর। বেল, নেবো ও শামাশ রাজা, আমার প্রভুর প্রতি সদয় হোন। একটি গ্রহণ ঘটেছে, কিন্তু রাজধানীতে তা দেখা যায়নি। সেই গ্রহণ যখন এগিয়ে আসছিল, তখন রাজা যেখানে বাস করেন সেই রাজধানীতে, দেখুন, সর্বত্র মেঘ ছিল। গ্রহণ ঘটেছে কি ঘটেনি তা আমরা জানি না। রাজাদের প্রভু আশুরে, সব নগরে, ব্যাবিলন, নিপ্পুর, উরুক ও বোরসিপ্পায় লোক পাঠান; সেই নগরগুলোতে যা দেখা গেছে রাজা তা নিশ্চিতভাবে শুনবেন।… রাজা, আমার প্রভুর নগরে বসবাসকারী মহান দেবতারা আকাশ মেঘে ঢেকে দিয়েছেন এবং গ্রহণ দেখতে দেননি। তাই রাজা জানুন যে এই গ্রহণ রাজা, আমার প্রভু, বা তাঁর দেশের বিরুদ্ধে নয়। রাজা আনন্দ করুন।…

এটি কি চমৎকার একটি ভাবনা নয় যে গ্রহণ ঘটার সময় আশুরের রক্ষাকর্তা দেবতারা, দেশ ও রাজার কোনো বিপদ নেই জেনে, নগরের ওপর মেঘের পর্দা টেনে দিলেন, যাতে রাজা ভয় না পান? লাইডেনের অ্যাসিরিওলজিস্ট দে লিয়াগ্রে ব্যোলের মতে, এই প্রতিবেদনগুলোর বেশিরভাগই খ্রিস্টপূর্ব ৬৭৫ থেকে ৬৬৫ সালের মধ্যেকার। তখন ছিল আসিরীয় সাম্রাজ্যের সংকটকাল; অসুস্থ ও আশঙ্কাগ্রস্ত রাজা কমবেশি ব্যর্থ একটি মিশর অভিযান থেকে ফিরেছিলেন। এখানে আমরা একটি বেশ নির্ভুল তারিখ পাই, যে সময়ের জ্যোতির্বৈজ্ঞানিক জ্ঞানের অবস্থা বলা যায়। এটা সত্যিই জ্ঞান ছিল, কারণ এই প্রতিবেদনগুলোতে গ্রহণ সম্পর্কে ভবিষ্যদ্বাণী করা হয়েছে এবং খুব স্পষ্টভাবে তা যাচাইও করা হয়েছে।

রাজা, আমার প্রভুর কাছে আমি পাঠিয়েছিলাম: "একটি গ্রহণ ঘটবে।" এখন তা ব্যর্থ হয়নি, ঘটেছে। এই গ্রহণ ঘটায় এটি রাজা, আমার প্রভুর জন্য শান্তির পূর্বাভাস দেয়।…

পুরোহিতেরা কিসের ভিত্তিতে ভবিষ্যদ্বাণী করতেন, এই লেখাগুলোতে তার কোনো ইঙ্গিত নেই। তবে চন্দ্রগ্রহণ ঘটার মধ্যে এমন একটি সরল নিয়মিততা আছে (যা প্রথম দেখিয়েছিলেন স্কিয়াপারেল্লি) যে তা নিশ্চয়ই শিগগিরই তাঁদের নজরে পড়েছিল। তাঁরা দেখতে পেতেন, ছয় মাসের কম ব্যবধানে কখনো দ্বিতীয় গ্রহণ হয় না, আর প্রায়ই আরও ছয় মাস পর তৃতীয় একটি গ্রহণ হয়। কখনো কখনো এমনকি চার বা পাঁচটি গ্রহণ পরপর ঘটত, সবসময় ছয় মাসের ব্যবধানে। প্রায়ই এই ধারাবাহিক শ্রেণিতে একটি ফাঁক থাকত। এমন হতো যখন গ্রহণটি দিনের বেলায় ঘটত, ফলে পূর্ণিমার চাঁদ দিগন্তের নিচে থাকত। পর্যবেক্ষকেরা একবার যখন বুঝলেন যে এটাই ফাঁকের কারণ, তখন সঠিক মুহূর্তে একটি অদৃশ্য গ্রহণ যোগ করে তাঁরা ফাঁকটি পূরণ করতে পারতেন। এভাবে পাঁচ বা ছয়টি দৃশ্যমান বা অদৃশ্য গ্রহণের একটি নিয়মিত ও অবিচ্ছিন্ন শ্রেণি তৈরি হতো। তারপর শ্রেণিটি শেষ হয়ে যেত। কিন্তু এক বা দুই বছর পর একটি নতুন শ্রেণি শুরু হতো, আগের শ্রেণিটি চলতে থাকলে যে সময়ে হতো তার এক মাস আগে। তাই যদি তাঁরা এমন একটি গ্রহণ পর্যবেক্ষণ করতেন যার ৬, ১২ বা ১৮ মাস আগে আর কোনো গ্রহণ হয়নি, তাহলে সিদ্ধান্ত নিতে পারতেন যে একটি নতুন শ্রেণি শুরু হয়েছে। এভাবে ৬ বা ১২ মাস পরের নতুন গ্রহণের ভবিষ্যদ্বাণী করতে পারতেন।

আধুনিক জ্ঞানের আলোকে এই নিয়মিততার ভিত্তি সহজেই বোঝা যায়। ছয়টি চান্দ্র পর্যায়কাল (১৭৭.১৮ দিন) পর সূর্য, আর তাই পূর্ণিমার চাঁদের অবস্থানও, দ্রাঘিমায় গড়ে ১৭৪.৬৪° এগিয়ে যায়। সূর্য ও চাঁদের পরিবর্তনশীল বেগের কারণে প্রকৃত অগ্রগতি কয়েক ডিগ্রি কম বা বেশি হতে পারে। একই সময়ে চন্দ্রকক্ষের পাতবিন্দুগুলো (nodes) ৯.৩৮° পিছিয়ে যায়; তাই বিপরীত পাতবিন্দুটি ১৭০.৬২° বেশি দ্রাঘিমায় অবস্থান করে। পাতবিন্দুর সাপেক্ষে পূর্ণিমার চাঁদ তখন ৪.০২° এগিয়ে যায়। চাঁদ যদি কোনো পাতবিন্দুর কাছে থাকে, তাহলে ছয় মাস পর আবারও একটি পাতবিন্দুর কাছে থাকবে। পাতবিন্দু থেকে পূর্ণিমার চাঁদের দূরত্ব ১০° থেকে ১২°-এর বেশি না হলে চাঁদ পৃথিবীর ছায়ার মধ্য দিয়ে যেতে পারে। পূর্ণগ্রহণের জন্য দূরত্ব ৫° বা ৬°-এর বেশি হওয়া চলবে না, নইলে গ্রহণটি কেবল আংশিক হয়। ছয় মাসের ব্যবধানে পরপর কয়েকটি পূর্ণিমার দূরত্ব হিসাব করলে এই শর্তগুলো স্পষ্ট করা যায়। হিসাবটি যেকোনো একটি প্রথম মান থেকে শুরু করে প্রতিবার ৪.০২ বা ৪.০৩ করে বাড়িয়ে করা হয়:

দূরত্ব	−১৪.৫০	−১০.৪৮	−৬.৪৫	−২.৪৩	+১.৫৯	+৫.৬২	+৯.৬৫	+১৩.৬৭
ফল	গ্রহণ নেই	আংশিক?	আংশিক	পূর্ণ	পূর্ণ	পূর্ণ?	আংশিক	গ্রহণ নেই

সংশ্লিষ্ট চিত্রে (২) এই অবস্থাগুলো দেখানো হয়েছে। সেখানে বড় বৃত্তগুলো ছায়া এবং ছোট বৃত্তগুলো চাঁদকে বোঝায়।

\end{document}